\documentclass[review]{elsarticle}
% Compile with: pdflatex draft.tex (run twice for cross-references)
% "review" option = double-spaced, numbered lines, single column — standard for JSS submission.
% Swap to \documentclass[preprint,5p]{elsarticle} for camera-ready double-column layout.
%
% NOTE ON DRAFT STATUS (author-facing, not part of the manuscript):
% This is a fully reconciled draft with respect to internal consistency (simulator identity,
% seeds, thresholds, RQ numbering, gate semantics). Negative or mixed results are reported as
% such rather than adjusted to fit an earlier hypothesis: the stratified-correlation (S5.5, S8.2)
% and shared-library blast-radius (S5.4, S8.2) analyses did not confirm the effects they were
% designed to detect, and the remediation impact quantification (S6.4, S6.7) is a mixed result
% (mean +4.61% across scenarios, but negative in 3 of 7) exposing a real gap between the per-edit
% acceptance test S6.4 describes and what PrescribeService currently does. The external,
% real-world validation planned for a prior draft (an ICAO-compliant air-traffic-management case
% study with a blind expert-ranking panel) has been withdrawn from this submission: the panel was
% never convened, and previously reported figures were placeholders that must not be presented as
% an executed study. Genuine external validation is left to future work (S9.3).

\usepackage{amsmath}
\usepackage{amssymb}
\usepackage{hyperref}
\usepackage{booktabs}
\usepackage{array}

\journal{Journal of Systems and Software}

\begin{document}

\begin{frontmatter}

\title{Graph Neural Networks for Reliability and Dependability Analysis in Complex Distributed Systems based on Publish--Subscribe Architecture}

\author{Author Name(s) Here}
\address{Affiliation Here}

\begin{abstract}
Modern distributed systems increasingly rely on publish--subscribe middleware to decouple data
producers and consumers. While this decoupling provides scaling and operational flexibility, it
obscures the true dependency chains along which a single component's failure can cascade.
Identifying \emph{which} components are critical --- and \emph{why} --- before deployment is
difficult: runtime telemetry does not yet exist at design time, and code-level Static Code Analysis
(SCA) platforms (e.g., SonarQube) are blind to system-level topological dependencies.

To bridge this ``Architecture-Code Gap,'' we present \textbf{Software-as-a-Graph (SaG)}, a
pre-deployment \textbf{Static System Analysis (SSA)} framework that models a distributed pub-sub
system as a typed, weighted, directed multigraph over five component classes (applications,
libraries, topics, brokers, and deployment nodes) and derives logical \texttt{DEPENDS\_ON}
dependencies through a set of typed projection rules.

On this typed representation, we employ a \textbf{Heterogeneous Graph Neural Network (GNN)}
predictor --- a relation-specific Graph Transformer with explicit Quality-of-Service (QoS) contract
edge-feature injection ($HGL\text{-}QoS$) --- to forecast each component's cascading failure impact
$I(v)$ before a single line of the system is deployed. We pair this learned predictor with an
interpretable, Analytic Hierarchy Process (AHP)-weighted multi-dimensional quality attribution score
$Q(v)$ that decomposes criticality into orthogonal Reliability, Maintainability, Availability, and
Vulnerability (RMAV) dimensions, so that every diagnostic is traceable to a concrete remediation.
Both predictors are validated against an independent discrete-event cascade simulator under a strict
input--label independence guarantee, ruling out transductive leakage.

Evaluated across seven synthetic, industrially-styled pub-sub topologies, we show: (1)
in-distribution regimes are strongly modeled by the deterministic $Q(v)$ attribution ($\rho > 0.87$,
$F_1 > 0.90$), competitive with or exceeding the learned predictor's in-distribution mean;
out-of-distribution generalization to entirely unseen architectures instead requires typed graph
learning, where $HGL\text{-}QoS$ decisively outperforms homogeneous and non-typed baselines,
reaching a Leave-One-Scenario-Out (LOSO) rank correlation of $\rho = 0.401$ against $\rho \approx 0$
for homogeneous GATs; (2) a stratified correlation check confirms the predictive relationship
between $Q(v)$ and simulated impact holds at consistent, moderate strength across all five node
types (pooled $\rho = 0.374$; per-type $\rho = 0.322$--$0.429$), and a direct test of the
hypothesized shared-library ``simultaneous blast'' mismatch does not find it in this suite ---
both reported as findings rather than adjusted to fit prior expectations; and (3) a multi-seed,
end-to-end evaluation of topology-level prescriptive remediation operators yields a mixed result: a
cross-scenario mean impact reduction of $+4.61\%$ that conceals resilience regressions of up to
$-31.67\%$ in 3 of 7 scenarios, exposing a concrete gap between the per-edit acceptance test the
design specifies and the unconditional policy application the current implementation performs.

Finally, we demonstrate SaG's operational feasibility as a delta-aware, blocking CI/CD quality gate:
using a thread-safe, database-free \texttt{MemoryRepository} to eliminate live database latency, the
gate evaluates complex regressions in $\approx 5\,\text{s}$ (medium topologies) to $\approx
40\,\text{s}$ (hyper-scale topologies) while achieving 100\% precision and recall on injected
structural regressions, enabling continuous, automated dependability auditing.
\end{abstract}

\begin{highlights}
\item Typed multigraph model bridges the Architecture-Code Gap in pub-sub systems.
\item AHP-weighted RMAV attribution explains criticality along four orthogonal axes.
\item Heterogeneous Graph Transformer (HGL-QoS) wins where interpretable scoring cannot.
\item Delta-aware CI/CD gate blocks structural regressions in seconds, pre-deployment.
\item Negative and mixed results (blast mismatch, remediation gaps) reported honestly.
\end{highlights}

\begin{keyword}
publish--subscribe middleware \sep architectural dependability \sep cascading failure \sep
heterogeneous graph neural networks \sep static system analysis \sep pre-deployment verification
\sep quality attributes \sep CI/CD quality gate
\end{keyword}

\end{frontmatter}

% ===========================================================================
\section{Introduction}
% ===========================================================================

\subsection{Motivation}

The publish--subscribe (pub-sub) paradigm has become a backbone communication abstraction for
large-scale distributed systems, underpinning cyber-physical, cloud-native, robotics, and
Internet-of-Things architectures. Its appeal is decoupling: producers and consumers are separated in
time, space, and synchronization, so components can be added, removed, or scaled without direct
knowledge of one another \cite{Eugster2003}. Industry standards such as the Data Distribution
Service (DDS) and MQTT formalize this model and expose deployment-time choices --- topics, brokers,
reliability, durability, and other Quality-of-Service (QoS) policies --- that materially shape how
the system behaves under stress \cite{OMG2015DDS,OASIS2019MQTT}.

The same decoupling that makes pub-sub flexible also obscures the dependency structure an engineer
must reason about when a component fails. There are no explicit caller--callee edges: an application
that publishes to a topic has no static link to the applications that subscribe to it, even though
those subscribers are wholly dependent on it for data. Failures do not propagate along a call graph
but along \emph{derived} paths --- through shared topics and brokers, through colocated deployment
nodes, and, distinctively, through shared libraries whose failure strikes every consumer
\emph{simultaneously} rather than sequentially. A raw architecture diagram does not reveal these
chains, and the components whose failure would be most damaging are frequently not the ones a
diagram makes look important.

Crucially, the moment at which this reasoning is most valuable is \emph{before} deployment.
Architectural hardening --- replication, isolation, failover, additional monitoring --- is cheapest
and least disruptive while the system is still a design, and prohibitively expensive once it is in
production. Yet pre-deployment is precisely when no runtime telemetry exists to identify weak points
empirically. An engineer must therefore answer a hard question from the architecture alone:
\emph{which components are critical, and why?}

\subsection{The Architecture-Code Gap and Problem Statement}

We address pre-deployment criticality analysis for pub-sub middleware as two coupled sub-problems.
Given only an architectural description of the system --- its applications, libraries, topics,
brokers, deployment nodes, and the QoS policies on its communication --- we seek to:

\begin{enumerate}
\item \textbf{Quality attribution.} Assign each component an interpretable measure of \emph{how} and
  \emph{why} it is critical, decomposed along the quality dimensions an engineer would act on, so
  that the result directs a specific remediation rather than a generic warning.
\item \textbf{Failure-impact analysis.} Predict the cascade impact of each component's failure ---
  the extent to which the rest of the system becomes unreachable or impaired --- and identify the
  components that should be hardened first.
\end{enumerate}

Both must be computed without runtime data, and both must remain \emph{explainable}: a single opaque
criticality number is of limited use to an architect who has to choose between competing
interventions under a fixed budget.

Historically, static verification has operated primarily at the source-code level. However, a major
``\textbf{Architecture-Code Gap}'' exists: a software system can have perfectly clean source code in
every component (earning top scores on code-level tools), yet remain highly fragile. If the
deployment topology contains a Single Point of Failure (SPOF) or a mismatched QoS contract, a single
component crash can cascade and collapse the entire system. Bridging this gap requires shifting
structural verification ``left'' into the continuous integration and delivery (CI/CD) pipeline.

\subsection{Limitations of Existing Approaches}

Three strands of prior work bear on this problem, and each leaves a gap.

\textbf{Static Code Analysis (SCA).} Platforms such as SonarQube evaluate code cleanliness,
cyclomatic complexity, and LCOM (Lack of Cohesion of Methods) inside individual modules. While highly
effective for intra-component quality, they are entirely blind to inter-component topologies and
dynamic middleware cascades.

\textbf{Runtime Dependability and Chaos Engineering.} A large body of work hardens pub-sub systems
through runtime fault tolerance, replication, and chaos injection (e.g., Chaos Monkey). These
techniques are valuable but assume a \emph{running} staging or production system; they do not answer
which components a design should protect before it is deployed, and injecting failures at runtime
carries operational risk.

\textbf{Topology-Only and Homogeneous Learning-Based Centrality.} Classical network-science metrics
collapse a component's risk into a single scalar that conflates distinct failure mechanisms (SPOFs
vs. cascade hubs), while homogeneous graph neural networks collapse typed semantics (applications,
topics, brokers) into flattened views, leading to representation collapse.

No existing approach offers an \emph{interpretable, multi-dimensional, pre-deployment} attribution
over the \emph{typed} pub-sub graph, coupled to code-level SCA metrics, heterogeneous-GNN impact
prediction, and automated CI/CD gating. That is the gap this paper fills.

\subsection{Our Approach}

We present \textbf{Software-as-a-Graph (SaG)}, a pre-deployment \textbf{Static System Analysis
(SSA)} framework. SaG models a pub-sub system as a typed, weighted, directed multigraph over five
node types (applications, libraries, topics, brokers, nodes) and derives logical
\texttt{DEPENDS\_ON} dependencies through typed projection rules. Crucially, SaG ingests code-level
SCA metrics as vertex attributes and performs \textbf{multi-dimensional quality attribution},
decomposing criticality into orthogonal Reliability, Maintainability, Availability, and
Vulnerability (RMAV) dimensions using Analytic Hierarchy Process (AHP) weights.

SaG then performs \textbf{failure-impact analysis}, predicting cascade impact $I(v)$ with two
predictors: the AHP-weighted composite $Q(v)$ and a learned \textbf{Heterogeneous Graph Transformer}
($HGL\text{-}QoS$). Both are validated against a discrete-event simulator under an
\textbf{input--label independence guarantee}. Finally, a \textbf{prescriptive remediation} stage
generates topology-level hardening edits and verifies them on counterfactual graphs in-memory.

To make SSA continuous, SaG integrates directly into CI/CD pipelines as a \emph{delta-aware}
blocking gate. By utilizing a thread-safe, database-free \texttt{MemoryRepository} to bypass Neo4j
database overhead during build time, SaG executes anti-pattern scans and counterfactual simulations
in seconds, and fails the build (exit code 2) when a change \emph{introduces new, unwaived} CRITICAL
or HIGH severity structural anomalies relative to the merge base (\S\ref{sec:gate}).

Concretely, the paper is organized around four research questions:

\begin{quote}
\textbf{RQ1.} When does interpretable, multi-dimensional attribution suffice for pre-deployment
criticality prediction, and when is a learned heterogeneous graph neural network required to
recover the simulated impact ordering?

\textbf{RQ2.} What failure modes does multi-dimensional quality attribution expose that a
single-score topological centrality misses?

\textbf{RQ3.} How does explicit multi-attribute QoS contract feature injection affect
in-distribution convergence versus out-of-distribution Leave-One-Scenario-Out (LOSO)
generalizability?

\textbf{RQ4.} What is the feasibility and performance overhead of deploying the graph-based analyzer
as a blocking Quality Gate in continuous integration/delivery (CI/CD) pipelines?
\end{quote}

RQ1, RQ2, and RQ3 are answered on the synthetic scenario suite (\S\ref{sec:results}.1--\S\ref{sec:results}.3); RQ4 evaluates gating
feasibility and performance (\S\ref{sec:results}.4).

\subsection{Contributions}

This paper makes the following contributions:

\begin{enumerate}
\item \textbf{A typed graph model with hierarchical SCA metric integration.} We define the SaG
  multigraph and the RMAV decomposition, which propagates code-level quality metrics (SonarQube
  \texttt{cm\_*} fields) into global system criticality scores (\S\ref{sec:model}, \S\ref{sec:rmav}).
\item \textbf{An automated, delta-aware CI/CD Quality Gate.} We formulate a build-blocking gate that
  evaluates system-level risk statically, blocks only newly introduced and unwaived structural
  regressions relative to the merge base, and executes in seconds via an in-memory repository
  (\S\ref{sec:remediation}).
\item \textbf{Failure-impact analysis under an independence guarantee.} We predict cascade impact
  using both the interpretable $Q(v)$ and a learned Heterogeneous Graph Transformer
  ($HGL\text{-}QoS$), validated against discrete-event simulation (\S\ref{sec:impact}).
\item \textbf{A prescriptive remediation stage.} We formalize a Generate$\to$Verify procedure with
  four concrete operators and a per-edit acceptance test against counterfactual simulation
  (\S\ref{sec:acceptance}); we measure the current implementation end to end and report that the
  acceptance test is not yet wired in as a per-edit filter, which we show produces a mixed rather
  than uniformly positive result (\S\ref{sec:mixed-result}).
\item \textbf{A stratified evaluation of criticality correlation.} We report Spearman correlation
  between $Q(v)$ and simulated impact $I(v)$ separately by node type, and by pooled aggregate,
  showing the two agree in this suite (no divergence between pooled and per-type figures), which is
  itself a useful methodological check for evaluating heterogeneous architectures
  (\S\ref{sec:impact}, \S\ref{sec:results}).
\end{enumerate}

\subsection{Relationship to the Authors' Prior Work}

This work extends the authors' earlier structural baseline of the framework --- multi-layer graph
dependency analysis --- introduced in prior work \cite{AnonA}. The present paper consolidates that
structural foundation with the heterogeneous graph neural network predictor, multi-dimensional
quality attribution (\S\ref{sec:rmav}), SCA metric integration, and CI/CD gating and remediation
(\S\ref{sec:remediation}) into a single, self-contained submission targeted at this special issue's
focus on AI for reliability and dependability analysis; no companion manuscript reporting the
heterogeneous-GNN predictor is submitted or under review in parallel with this paper.

\subsection{Organization}

The remainder of this paper is organized as follows. Section~\ref{sec:relatedwork} reviews related
work. Section~\ref{sec:model} defines the Software-as-a-Graph model. Section~\ref{sec:rmav}
presents multi-dimensional quality attribution, and Section~\ref{sec:impact} presents
failure-impact analysis. Section~\ref{sec:remediation} introduces prescriptive remediation and
CI/CD quality gating. Section~\ref{sec:setup} describes the experimental setup;
Section~\ref{sec:results} reports the synthetic-suite and gating results (RQ1--RQ4); and
Section~\ref{sec:discussion} discusses the findings, threats to validity, and conclusions.

% ===========================================================================
\section{Related Work}
\label{sec:relatedwork}
% ===========================================================================

This paper draws on, and contributes to, several established lines of research: publish--subscribe
dependability, static analysis techniques, pre-deployment system verification, structural
criticality, and multi-criteria quality scoring.

\subsection{Publish--Subscribe Middleware and Dependability}

The pub-sub paradigm is a foundational communication abstraction for large-scale distributed
systems, valued for decoupling producers and consumers in time, space, and synchronization
\cite{Eugster2003}. Content-based and brokered overlays extend this with flexible event routing and
subscription matching, and standards such as DDS and MQTT formalize deployment-time choices ---
topics, brokers, reliability, durability, and other QoS policies --- that govern runtime behavior
\cite{OMG2015DDS,OASIS2019MQTT}. These mechanisms enable cyber-physical, cloud, IoT, and robotics
architectures, but they also make failure propagation difficult to reason about from direct
communication edges alone.

Research on pub-sub dependability has accordingly emphasized runtime fault tolerance, reliable event
dissemination, replication, and recovery. These approaches improve a system's resilience while it is
\emph{running}: they assume observable behavior and react to or mask faults as they occur. Our
concern is complementary and earlier in the lifecycle --- estimating, from an architectural model
that enumerates applications, libraries, topics, brokers, and QoS policies, which components would
have the greatest downstream impact if they failed, so that the design can be hardened before any
system is deployed.

\subsection{Static Code Analysis (SCA) vs. Static System Analysis (SSA)}

Static verification typically operates at the source-code level. Static Code Analysis (SCA) tools,
exemplified by SonarQube, checkstyle, and FindBugs, parse source files into Abstract Syntax Trees
(ASTs) to compute complexity, code duplication, and modular metrics such as LCOM (Lack of Cohesion
of Methods). While SCA is essential for locating intra-component defects and technical debt, it is
blind to the inter-component topology.

Static System Analysis (SSA) addresses this ``Architecture-Code Gap.'' SSA models the system as a
global graph of communicating components, middleware routers, and hardware hosts. Rather than
replacing SCA, SSA ingests code-level metrics as node properties (e.g., LCOM, cyclomatic complexity)
and propagates them through the inter-component dependency topology. This allows architects to
evaluate how code-level fragility (e.g., a highly complex class inside an application) combines with
structural fragility (e.g., the application being a single point of failure) to create systemic
risks.

\subsection{Continuous Pre-Deployment Verification and Gating}

A common way to verify system resilience is dynamic testing, particularly Chaos Engineering (e.g.,
Netflix Chaos Monkey), which injects faults into live staging or production clusters. While chaos
testing evaluates real operational environments, doing so carries risk and occurs late in the
lifecycle.

Continuous pre-deployment verification shifts this analysis left, integrating it into CI/CD
pipelines (e.g., GitHub Actions, GitLab CI). In this paradigm, the system architecture is defined as
``Architecture-as-Code'' (AaC) via configuration descriptors (Docker Compose, Kubernetes manifests,
Helm charts). SSA tools run automatically on every pull request, parsing the configuration
descriptors to generate a counterfactual topology graph, and block the build (exiting with non-zero
status) when a change introduces critical architectural smells (like SPOFs or QoS mismatches) or
exceeds failure-propagation thresholds. Mature code-level gates follow the same discipline:
SonarQube's default ``Clean as You Code'' quality gate evaluates \emph{new} code against the merge
base rather than failing builds on the accumulated state of the whole codebase, and pairs the gate
with an explicit won't-fix/false-positive marking workflow. Our system-level gate adopts the
analogous semantics --- blocking on newly introduced, unwaived structural regressions rather than on
any pre-existing finding (\S\ref{sec:gate}) --- because real architectures legitimately contain
\emph{intentional}, risk-accepted SPOFs that an absolute gate would flag on every build.

\subsection{Structural Criticality Analysis}

Network science offers a mature toolkit for identifying important nodes and edges. Degree, closeness
and betweenness centrality, articulation points, and PageRank-style scores are prized for their
efficiency and interpretability \cite{Freeman1977,BrinPage1998}, and studies of node removal,
cascading failure, and interdependent networks have deepened our understanding of systemic fragility
\cite{Buldyrev2010}. Applied to software dependency graphs, these metrics can flag bottlenecks and
single points of failure at design time.

Their limitation, for our purpose, is dimensional collapse. A single centrality score conflates
mechanisms that call for different remedies: a structural single point of failure, a high-reach
cascade hub, and a tightly coupled maintainability bottleneck can all present as ``central,'' yet a
replica, a rerouting, and a decoupling refactor are not interchangeable fixes. Untyped metrics are
moreover blind to type-specific failure modes --- most importantly the \emph{simultaneous} blast
radius of a shared library, whose failure strikes all consumers at once and which is
indistinguishable from an ordinary edge once node and edge types are discarded. Our RMAV attribution
retains the interpretability that makes structural metrics attractive while decomposing criticality
into orthogonal dimensions, and our typed model keeps the semantics that single-score centrality
erases.

\subsection{Learning-Based Criticality Prediction}

A growing body of work learns to identify critical nodes directly from graph structure, often
surpassing hand-crafted metrics when higher-order structure matters: FINDER locates key entities in
networked systems, DrBC learns to approximate betweenness, and PowerGraph provides a GNN benchmark
for cascading-failure and critical-node analysis in power-grid networks
\cite{Fan2020Nature,Fan2019CIKM,Varbella2024PowerGraph}.

Most such methods, however, target \emph{homogeneous} graphs. Pub-sub middleware is intrinsically
heterogeneous --- applications publish and subscribe to topics, topics are routed through brokers,
libraries introduce code dependencies, and deployment nodes impose locality --- and flattening this
into a homogeneous graph discards information about how failures propagate. Heterogeneous graph
neural networks address this directly: RGCN applies relation-specific transformations
\cite{Schlichtkrull2018RGCN}, HAN uses hierarchical attention \cite{Wang2019HAN}, HGT parameterizes
attention by node and edge type \cite{Hu2020HGT}, and MAGNN aggregates along metapaths
\cite{Fu2020MAGNN}. A known hazard in dense, hub-dominated regions is over-smoothing
\cite{Li2018DeeperInsights}. Our learned predictor adopts relation-specific message passing over the
native typed architecture for exactly these reasons, but we treat it as one of two predictors rather
than the sole contribution: a central question of this paper (RQ1) is \emph{when} such learning is
necessary at all, given an interpretable alternative.

\subsection{Quality Attributes and Multi-Criteria Scoring}

Software quality is conventionally described along attributes such as reliability, maintainability,
availability, and security, and a substantial literature connects these attributes to measurable
structural and code-level properties. Combining several such properties into a single decision score
is a multi-criteria decision problem, for which the Analytic Hierarchy Process (AHP) provides a
principled, auditable weighting derived from pairwise comparisons with an explicit consistency check
\cite{Saaty1980}.

What has not been done, to our knowledge, is to use a multi-criteria decomposition as the
\emph{attribution} mechanism for pre-deployment component criticality in pub-sub systems --- that
is, to make the per-dimension breakdown the explanation an architect acts on, with each structural
metric feeding exactly one dimension so that the reason a component is critical is legible from its
profile. Our RMAV scoring does precisely this, using AHP both within each dimension and to form the
composite $Q(v)$, with a shrinkage toward a uniform prior to guard against extreme weights on small
comparison sets. This connects the interpretability tradition of structural analysis (\S2.4) to the
decision-theoretic tradition of multi-criteria scoring, and is what distinguishes attribution here
from an opaque learned score.

\subsection{Architectural Remediation and Anti-Pattern Detection}

A related strand detects architectural anti-patterns and recommends refactorings --- cyclic
dependencies, hubs, unstable interfaces --- typically from a static dependency model, and evaluates
the effect of a change by re-analyzing the modified model. Our prescriptive stage is in this spirit
but differs in its acceptance test: rather than accepting an edit because it improves a static
metric, we \emph{verify} each candidate edit on a counterfactual graph using the same discrete-event
simulation oracle that produces our ground-truth impact, and accept it only if the reduction in
simulated impact exceeds a multi-seed variance threshold. Generation of candidate edits remains
topology-only, preserving the independence between the diagnostic and validation paths that the rest
of the framework relies on.

\subsection{Positioning}

In summary, prior approaches either (i) address pub-sub dependability at the protocol or runtime
level, presupposing a deployed system; (ii) offer code-level SCA that is blind to inter-component
topologies; (iii) offer structural analysis that conflates failure mechanisms and ignores typed
modes such as shared-library blast; (iv) apply graph learning while discarding the typed semantics
of pub-sub; or (v) use multi-criteria scoring for prioritization but not as an interpretable
criticality \emph{attribution} over a typed architecture graph. Software-as-a-Graph combines a typed
multigraph model, AHP-based multi-dimensional attribution, dual interpretable and learned impact
predictors, and a simulation-verified, delta-aware continuous CI/CD quality gate. The stratified
correlation evaluation we report --- by node type as well as pooled --- is a direct consequence of
taking node and edge type seriously, and is a methodological standard the untyped or
single-dimensional methods reviewed above do not apply.

% ===========================================================================
\section{The Software-as-a-Graph Model}
\label{sec:model}
% ===========================================================================

This section defines the graph model on which all subsequent analysis operates. We first give the
formal object and its node and edge types (\S\ref{sec:model}.1), then the QoS-derived edge and
vertex weights that encode coupling strength (\S\ref{sec:model}.2), then the derivation of logical
dependencies from structural edges (\S\ref{sec:model}.3), the ingestion of code-level SCA metrics
(\S\ref{sec:model}.4), and finally the two graph views and the multi-layer projections that the
attribution and impact stages consume (\S\ref{sec:model}.5). A running example threads through the
section (\S\ref{sec:model}.6).

\subsection{Nodes, Edges, and the Formal Object}

A distributed publish--subscribe system is modeled as a typed, weighted, directed multigraph
\begin{equation}
G = (V, E, \tau_V, \tau_E, w_E, w_V),
\end{equation}
where the vertex set partitions into five component types,
\begin{equation}
V = V_{\text{app}} \cup V_{\text{broker}} \cup V_{\text{topic}} \cup V_{\text{node}} \cup V_{\text{lib}},
\end{equation}
the type functions $\tau_V : V \to \{\text{App}, \text{Broker}, \text{Topic}, \text{Node},
\text{Library}\}$ and $\tau_E$ label vertices and edges, and the weight functions $w_E : E \to
[0,1]$ and $w_V : V \to [0,1]$ encode QoS-derived coupling strength. The edge set is the disjoint
union of \emph{structural} edges imported directly from the architecture description and
\emph{dependency} edges (\texttt{DEPENDS\_ON}) derived from them (\S\ref{sec:model}.3).

\textbf{Node types.} Each type corresponds to a distinct architectural element with its own failure
semantics (Table~\ref{tab:nodetypes}).

\begin{table}[h]
\centering
\caption{Node types and their failure semantics.}
\label{tab:nodetypes}
\begin{tabular}{@{}p{2.2cm}p{5cm}p{5cm}@{}}
\toprule
Type & Role & Representative instances \\
\midrule
Application & A process that publishes and/or subscribes to topics & ROS~2 node, Kafka producer/consumer, MQTT client \\
Broker & A message-routing intermediary & RabbitMQ, Mosquitto, DDS middleware \\
Topic & A named message channel & \texttt{/sensor/lidar}, \texttt{order.events} \\
Node & A physical or virtual host & server, cloud VM, embedded controller \\
Library & A shared code dependency & sensor driver, codec, message library \\
\bottomrule
\end{tabular}
\end{table}

\textbf{Structural edge types.} Six edge types are imported from the topology description and carry
the direction in which messages or hosting relationships flow (Table~\ref{tab:edgetypes}).

\begin{table}[h]
\centering
\caption{Structural edge types.}
\label{tab:edgetypes}
\begin{tabular}{@{}p{3.2cm}p{2.8cm}p{6cm}@{}}
\toprule
Edge & Direction & Meaning \\
\midrule
\texttt{PUBLISHES\_TO} & App/Library $\to$ Topic & component produces messages on the topic \\
\texttt{SUBSCRIBES\_TO} & App/Library $\to$ Topic & component consumes messages from the topic \\
\texttt{ROUTES} & Broker $\to$ Topic & broker routes the topic \\
\texttt{RUNS\_ON} & App/Broker $\to$ Node & component is hosted on the node \\
\texttt{CONNECTS\_TO} & Node $\to$ Node & direct network link between hosts \\
\texttt{USES} & App $\to$ Library & application depends on the shared library \\
\bottomrule
\end{tabular}
\end{table}

Retaining these types --- rather than collapsing them into a single ``communicates-with'' relation
--- is what later lets the framework distinguish failure mechanisms that an untyped graph cannot
(\S\ref{sec:model}.3, \S\ref{sec:impact}).

\subsection{QoS-Aware Edge and Vertex Weights}

Not all dependencies are equally consequential: a \texttt{RELIABLE}/\texttt{PERSISTENT} channel
carrying critical data couples its endpoints far more tightly than a
\texttt{BEST\_EFFORT}/\texttt{VOLATILE} one. Edge weights encode this from the Quality-of-Service
policy of each pub-sub relationship, via a two-stage computation:
\begin{align}
\text{QoS\_score} &= 0.30\,r + 0.40\,d + 0.30\,p, \\
\text{size\_norm} &= \min\!\left(\frac{\log_2(1 + \text{size\_kb})}{50},\ 1.0\right), \\
w(e) &= \beta\cdot\text{QoS\_score} + (1-\beta)\cdot\text{size\_norm}, \qquad \beta = 0.85,
\end{align}
where $r, d, p$ are the reliability, durability, and transport-priority scores of the mediating
topic, mapped from symbolic QoS values (Table~\ref{tab:qos}).

\begin{table}[h]
\centering
\caption{Symbolic QoS value mapping.}
\label{tab:qos}
\begin{tabular}{@{}p{3cm}p{9cm}@{}}
\toprule
Dimension & Symbolic value $\to$ score \\
\midrule
Reliability $r$ & \texttt{RELIABLE} $\to$ 1.0; \texttt{BEST\_EFFORT} $\to$ 0.0 \\
Durability $d$ & \texttt{PERSISTENT} $\to$ 1.0; \texttt{TRANSIENT} $\to$ 0.6; \texttt{TRANSIENT\_LOCAL} $\to$ 0.5; \texttt{VOLATILE} $\to$ 0.0 \\
Priority $p$ & \texttt{URGENT}/\texttt{CRITICAL}/\texttt{HIGHEST} $\to$ 1.0; \texttt{HIGH} $\to$ 0.66; \texttt{MEDIUM} $\to$ 0.33; \texttt{LOW} $\to$ 0.0 \\
\bottomrule
\end{tabular}
\end{table}

The intra-QoS sub-weights are AHP-derived: durability (0.40) outweighs reliability and priority
(0.30 each) because durability governs message-state survival --- the precondition for resilience
--- whereas reliability and priority govern transient delivery quality. A floor of $w(e) = 0.01$
keeps even zero-QoS components visible to attribution.

\textbf{Vertex weights} propagate QoS upward from incident edges, with type-specific aggregation
that reflects how each component type concentrates risk:

\begin{table}[h]
\centering
\caption{Vertex weight aggregation by type.}
\label{tab:vertexweights}
\begin{tabular}{@{}p{2.2cm}p{9.5cm}@{}}
\toprule
Type & $w_V$ \\
\midrule
Application & $0.80\cdot\max(w_{\text{topic}}) + 0.20\cdot\operatorname{mean}(w_{\text{topic}})$ \\
Broker & $0.70\cdot\max(w_{\text{topic}}) + 0.30\cdot\operatorname{mean}(w_{\text{topic}})$ \\
Node & $\max(w)$ over all hosted applications and brokers \\
Library & $\min\!\big(1.0,\ w_{\text{base}}\cdot(1 + \gamma\log_2(1 + \mathrm{DG\_in}))\big)$ (fan-out amplified) \\
\bottomrule
\end{tabular}
\end{table}

The library rule is deliberately fan-out amplified: a library's risk grows with the number of
applications that depend on it, anticipating the blast-radius mechanism of \S\ref{sec:model}.3 and
\S\ref{sec:impact}.

\subsection{Derived Dependencies: the \texttt{DEPENDS\_ON} Projection}

Structural edges record physical relationships but not \emph{logical} dependency. A subscriber and a
publisher on the same topic have no direct structural edge, yet the subscriber wholly depends on the
publisher for data. We therefore derive a single semantic relation, \texttt{DEPENDS\_ON}, always
directed from \emph{dependent} to \emph{dependency} (``if the target fails, the source is
affected''), through six rules (Table~\ref{tab:rules}).

\begin{table}[h]
\centering
\caption{The six \texttt{DEPENDS\_ON} projection rules.}
\label{tab:rules}
\begin{tabular}{@{}cp{2.3cm}p{5.2cm}p{3cm}@{}}
\toprule
Rule & \texttt{dependency\_type} & Pattern (dependent $\to$ dependency) & Weight \\
\midrule
1 & \texttt{app\_to\_app} & subscriber $\to$ publisher via a shared topic (incl.\ transitive \texttt{USES*1..3} chains) & $\max_t w(t)$ over shared topics \\
2 & \texttt{app\_to\_broker} & publisher/subscriber $\to$ broker routing its topics & $\max_t w(t)$ over routed topics \\
3 & \texttt{node\_to\_node} & host $\to$ host, lifted from Rules 1--2 for colocated apps & lifted $\max w$ \\
4 & \texttt{node\_to\_broker} & host $\to$ broker, lifted from Rule 2 & lifted $\max w$ \\
5 & \texttt{app\_to\_lib} & application $\to$ library it \texttt{USES} --- \textbf{shared-library blast} & $w_V(\text{app})$ \\
6 & \texttt{broker\_to\_broker} & bidirectional, two brokers sharing a host --- \textbf{colocation} & $w_V(\text{node})$ \\
\bottomrule
\end{tabular}
\end{table}

When two applications communicate over several shared topics, a single \texttt{DEPENDS\_ON} edge
records the worst-case weight together with a separate coupling count:
\begin{equation}
\text{edge.weight} = \max_{t \in \text{shared}} w(t), \qquad \text{edge.path\_count} = |\text{shared}|.
\end{equation}
\texttt{path\_count} is kept out of the weight to preserve the $w \in [0,1]$ contract; a
\texttt{path\_count} of 3 denotes three simultaneous failure vectors between the same pair, which is
structurally more fragile than three independent single-topic links.

\textbf{Two qualitatively different failure modes.} This is the crux of the model. Rule 1 encodes
\emph{sequential cascade}: a publisher's failure starves its subscribers, whose failure may in turn
affect their dependents, propagating step by step through topics and brokers. Rule 5 encodes a
\emph{simultaneous blast}: when a shared library fails, every application that uses it fails at
once, in a single event, not along a propagation path. An untyped graph cannot tell these apart ---
both look like ordinary edges --- yet they demand different predictions and different remedies.
Preserving the \texttt{app\_to\_lib} type (Rule 5) is what lets the framework represent this
simultaneous-blast mechanism at all, just as preserving \texttt{broker\_to\_broker} (Rule 6) makes
broker-colocation risk representable.

\subsection{Ingestion of Code-Level SCA Metrics}

To bridge the ``Architecture-Code Gap,'' SaG does not operate in isolation from source code.
Instead, the framework integrates code-level quality attributes directly into the graph model.
During the model-import stage, SaG queries static code analysis (SCA) APIs (e.g., SonarQube's web
API) or parses local SCA report artifacts to extract modular metrics for executable
\texttt{Application} and shared \texttt{Library} components.

These metrics are stored as flat properties prefixed with \texttt{cm\_*} on each component node:
\begin{itemize}
\item \texttt{cm\_total\_loc}: Total lines of code as reported by static analysis, providing a scale
  proxy.
\item \texttt{cm\_avg\_wmc}: Average Weighted Methods per Class, representing cognitive complexity.
\item \texttt{cm\_avg\_lcom}: Lack of Cohesion of Methods (on a raw $[0,1]$ scale), indicating how
  fragmented classes are.
\item \texttt{cm\_avg\_cbo}: Coupling Between Objects, indicating intra-component code coupling.
\item \texttt{cm\_avg\_rfc}: Response for a Class, measuring the number of methods invoked by a
  class.
\item \texttt{sqale\_debt\_ratio}: Technical debt ratio as a percentage of estimated rewrite time.
\item \texttt{bugs}: Count of static bugs identified in code.
\item \texttt{vulnerabilities}: Count of code-level security issues.
\end{itemize}

These properties are normalized across the component population during structural analysis
(\S\ref{sec:rmav}.2) and feed the \textbf{Code Quality Penalty (CQP)}, ensuring that local code
defects are mathematically combined with global structural dependencies.

\subsection{Graph Views and Multi-Layer Projections}

The construction produces \textbf{two complementary views} of the same system, and the separation
between them is load-bearing for the framework's validity:

\begin{itemize}
\item $G_{\text{structural}}$ --- the imported structural graph, used by the discrete-event
  simulator to generate ground-truth impact $I(v)$ (\S\ref{sec:impact}).
\item $G_{\text{analysis}}(\ell)$ --- the layer-projected \texttt{DEPENDS\_ON} graph, on which all
  structural metrics, quality attribution, and prediction are computed (\S\ref{sec:rmav}).
\end{itemize}

Because attribution is computed on $G_{\text{analysis}}$ while ground truth is generated by
simulating $G_{\text{structural}}$, the predictor's inputs are kept disjoint from the
label-producing path --- the \textbf{independence guarantee} that makes the pre-deployment claims of
\S\ref{sec:rmav}--\S\ref{sec:impact} non-circular. We state and rely on this property throughout.

$G_{\text{analysis}}$ is filtered into four analytical layers, each isolating a component scope, a
dependency subset, and the quality dimension it most informs (Table~\ref{tab:layers}).

\begin{table}[h]
\centering
\caption{Analytical layer projections.}
\label{tab:layers}
\begin{tabular}{@{}p{2.3cm}p{1.6cm}p{2.6cm}p{4cm}p{2cm}@{}}
\toprule
Layer & Projection & Vertices & Dependency types & Quality focus \\
\midrule
Application & $\pi_{\text{app}}$ & App, Library & \texttt{app\_to\_app}, \texttt{app\_to\_lib} & Reliability \\
Infrastructure & $\pi_{\text{infra}}$ & Node & \texttt{node\_to\_node} & Availability \\
Middleware & $\pi_{\text{mw}}$ & Broker (in App/Node context) & \texttt{app\_to\_broker}, \texttt{node\_to\_broker}, \texttt{broker\_to\_broker} & Maintainability \\
System & $\pi_{\text{system}}$ & all five types & all six & Overall \\
\bottomrule
\end{tabular}
\end{table}

The middleware layer includes Application and Node vertices in the subgraph to preserve incoming
edges, but reports results only for Brokers. Components further aggregate along a MIL-STD-498
hierarchy --- CSU $\to$ CSC $\to$ CSCI $\to$ CSS --- so that criticality can be rolled up from a unit
to a configuration item to the whole system, supporting reporting at whatever granularity an
organization's software configuration management already uses.

\subsection{Running Example}
\label{sec:runningexample}

Consider three applications $a_1, a_2, a_3$, where $a_1$ publishes to a topic $t$ that $a_2$ and
$a_3$ subscribe to, all three depending on a shared library $\ell$. The structural graph records
$a_1\xrightarrow{\text{pub}}t$, $a_2,a_3\xrightarrow{\text{sub}}t$, and
$a_i\xrightarrow{\text{uses}}\ell$. Derivation adds $a_2\to a_1$ and $a_3\to a_1$ (\texttt{app\_to\_app},
Rule~1) and $a_i\to\ell$ (\texttt{app\_to\_lib}, Rule~5). The two structures encode different risks:
losing $a_1$ degrades $a_2$ and $a_3$ through a cascade that the simulator propagates over time,
whereas losing $\ell$ fails $a_1, a_2, a_3$ simultaneously. A topology-only centrality score would
rank $\ell$ by ordinary connectivity and miss that its single failure collapses the whole component
group at once --- the gap \S\ref{sec:impact} quantifies.

% Figure placeholder: no image asset currently exists for this running example.
% \begin{figure}[h]
% \centering
% \includegraphics[width=0.8\textwidth]{running-example.pdf}
% \caption{Structural graph and its derived \texttt{DEPENDS\_ON} projection, with cascade and blast
% edges visually distinguished.}
% \label{fig:runningexample}
% \end{figure}

% ===========================================================================
\section{Multi-Dimensional Quality Attribution (The Interpretable Path)}
\label{sec:rmav}
% ===========================================================================

Centrality answers \emph{whether} a component is important with a single number. An architect
choosing between a replica, a reroute, and a decoupling refactor needs to know \emph{why}. This
section presents the framework's primary diagnostic: a decomposition of each component's
criticality into four orthogonal quality dimensions, each computed from disjoint structural metrics,
and combined into an interpretable composite score. Because the dimensions do not share inputs, a
component's profile is itself the explanation of its risk --- and the explanation maps directly to
a remedy (\S\ref{sec:remediation}).

\subsection{Four Orthogonal Dimensions}

We attribute criticality along Reliability, Maintainability, Availability, and Vulnerability (RMAV).
Each answers a distinct architectural question and speaks to a distinct stakeholder
(Table~\ref{tab:rmavdims}).

\begin{table}[h]
\centering
\caption{The four RMAV dimensions.}
\label{tab:rmavdims}
\begin{tabular}{@{}cp{3.4cm}p{3.6cm}p{2.6cm}@{}}
\toprule
Dim. & Question & High score means & Stakeholder \\
\midrule
\textbf{R} & How broadly and deeply does failure propagate? & Failure cascades widely; hard to contain & Reliability engineer \\
\textbf{M} & How hard is this to change safely? & Tightly coupled structural bottleneck & Software architect \\
\textbf{A} & Is this a structural single point of failure? & Removing it partitions the dependency graph & DevOps / SRE \\
\textbf{V} & How attractive a target is this for attack? & Central, reachable, high-value downstream & Security engineer \\
\bottomrule
\end{tabular}
\end{table}

The dimensions are \textbf{orthogonal by construction}: each raw structural metric feeds exactly one
dimension, never more. This is a deliberate design constraint, not an empirical observation ---
allowing a metric into two dimensions would silently inflate its weight relative to the AHP
calibration (\S\ref{sec:rmav}.3). Orthogonality is what makes the breakdown legible: a pure single
point of failure scores high on A but low on R, M, and V; a god-component scores high on M; a
cascade hub scores high on R. The \emph{shape} of the profile names the failure mode.

\subsection{RMAV Formulas}

All metric inputs are rank-normalized to $[0,1]$, so every RMAV score lies in $[0,1]$.
Table~\ref{tab:notation} fixes notation for every structural metric the four formulas below
consume; each is computed once on $G_{\text{analysis}}$ and feeds exactly one RMAV dimension
(\S\ref{sec:rmav}.1).

\begin{table}[p]
\centering
\scriptsize
\caption{RMAV input metric notation. $G^\top$ denotes the transpose of the \texttt{DEPENDS\_ON}
graph (the failure-propagation direction, since edges point dependent $\to$ dependency).}
\label{tab:notation}
\begin{tabular}{@{}p{2.0cm}p{2.6cm}p{5.2cm}p{0.7cm}@{}}
\toprule
Symbol & Name & Computed as & Feeds \\
\midrule
$\mathrm{RPR}(v)$ & Reverse PageRank & PageRank on $G^\top$ ($d=0.85$) & $R$ \\
$\mathrm{DG\_in}(v)$ & In-degree (rank-norm.) & Direct dependent count on \texttt{DEPENDS\_ON} & $R$ \\
$\mathrm{MPCI}(v)$ & Multi-Path Coupling Index & $\sum_{e\in\text{InEdges}(v)} \max(\text{path\_count}(e)-1,0) / (|V|-1)$ & $R$ \\
$\mathrm{CDPot\_enh}(v)$ & Enhanced Cascade Depth Potential & RPR/DG\_in blend, amplified by MPCI & $R$ \\
$\mathrm{FOC}(v)$ & Fan-Out Criticality & frequency- and QoS-weighted subscriber fan-out (Topic nodes only) & $R_{\text{topic}}$ \\
$\mathrm{BT}(v)$ & Betweenness centrality & Brandes' algorithm on $G_{\text{analysis}}$, QoS-inverted edge distances & $M$ \\
$w_{\text{out}}(v)$ & QoS-weighted out-degree & $\sum_{(v,u)} w(v,u)$ over outgoing dependencies & $M$ \\
$\mathrm{CQP}(v)$ & Code Quality Penalty & SonarQube-derived composite (\S\ref{sec:model}.4); 0 for non-App/Library types & $M$ \\
$\mathrm{CouplingRisk\_enh}(v)$ & Enhanced coupling risk & in/out-degree balance amplified by path complexity & $M$ \\
$\mathrm{CC}(v)$ & Clustering coefficient & Watts--Strogatz local clustering on the undirected projection & $M$ (as $1-\mathrm{CC}$) \\
$\mathrm{AP\_c\_directed}(v)$ & Directed articulation score & $\max$ of directed in/out articulation scores & $A$ \\
$\mathrm{QSPOF}(v)$ & QoS-weighted SPOF severity & $\mathrm{AP\_c\_directed}(v)\cdot w(v)$ & $A$ \\
$\mathrm{BR}(v)$ & Bridge ratio & fraction of $v$'s undirected edges that are bridges & $A$ \\
$\mathrm{CDI}(v)$ & Connectivity Degradation Index & normalized increase in average path length when $v$ is removed & $A$ \\
$\mathrm{REV}(v)$ & Reverse eigenvector centrality & eigenvector centrality on $G^\top$ & $V$ \\
$\mathrm{RCL}(v)$ & Reverse closeness (harmonic) & harmonic centrality on $G^\top$, normalized by $|V|-1$ & $V$ \\
$w_{\text{in}}(v)$ & QoS-weighted in-degree (QADS) & $\sum_{(u,v)} w(u,v)$ over incoming dependencies & $V$ \\
\bottomrule
\end{tabular}
\end{table}

\textbf{Reliability} --- fault-propagation risk. Because \texttt{DEPENDS\_ON} points \emph{dependent
$\to$ dependency}, a failure propagates \emph{against} edge direction; RPR (computed on the
transpose $G^\top$) therefore traverses the natural failure-propagation path. For Topic nodes, which
have no \texttt{DEPENDS\_ON} in-degree, a fan-out form is dispatched by $\tau_V(v)$:
\begin{align}
R(v) &= 0.45\cdot\mathrm{RPR}(v) + 0.30\cdot\mathrm{DG\_in}(v) + 0.25\cdot\mathrm{CDPot\_enh}(v)
\qquad [\tau_V(v)\neq\text{Topic}] \\
\mathrm{CDPot\_enh}(v) &= \min\!\Big( \frac{\mathrm{RPR}(v) + \mathrm{DG\_in}(v)}{2} \cdot \big(1 -
\min(\tfrac{\mathrm{out\_degree\_raw}(v)}{\max(\mathrm{in\_degree\_raw}(v),\, \epsilon)}, 1)\big)
\cdot (1 + \mathrm{MPCI}(v)),\ 1.0 \Big) \\
R_{\text{topic}}(v) &= 0.50\cdot\mathrm{FOC}(v) + 0.50\cdot\mathrm{CDPot\_topic}(v),\quad
\mathrm{CDPot\_topic}(v) = \mathrm{FOC}(v)\big(1 - \min(\text{publisher\_count\_norm}(v),1)\big)
\end{align}

\textbf{Maintainability} --- coupling complexity:
\begin{align}
M(v) &= 0.35\,\mathrm{BT}(v) + 0.30\,w_{\text{out}}(v) + 0.15\,\mathrm{CQP}(v)
+ 0.12\,\mathrm{CouplingRisk\_enh}(v) + 0.08\,(1-\mathrm{CC}(v)), \\
\mathrm{CQP}(v) &= 0.10\,\text{loc\_norm} + 0.35\,\text{complexity\_norm}
+ 0.30\,\text{instability\_code} + 0.25\,\text{lcom\_norm}.
\end{align}

Here, the Code Quality Penalty (CQP) translates local code-level fragility into system-level
maintainability risk. The components \texttt{loc\_norm}, \texttt{complexity\_norm}, and
\texttt{lcom\_norm} represent the min-max normalized values of the ingested SonarQube properties
\texttt{loc}, \texttt{cyclomatic\_complexity}, and \texttt{lcom}, respectively. These are calculated
independently for Applications and Libraries to prevent scale differences from distorting the
normalization. The metric \texttt{instability\_code} represents class instability (efferent coupling
divided by total coupling). The CQP thus ensures that local code debt is penalised, but only as a
sub-factor of Maintainability ($M$), which remains heavily weighted by topological metrics such as
betweenness centrality ($BT$) and efferent QoS-weighted out-degree ($w_{\text{out}}$). CQP is zero
for non-Application/Library types (graceful degradation). The two instability signals are
intentional and distinct: \texttt{instability\_code} is static-code fragility (local);
\texttt{CouplingRisk\_enh} is runtime-topology fragility (global).

\textbf{Availability} --- single-point-of-failure risk:
\begin{equation}
A(v) = 0.35\,\mathrm{AP\_c\_directed}(v) + 0.25\,\mathrm{QSPOF}(v) + 0.25\,\mathrm{BR}(v)
+ 0.10\,\mathrm{CDI}(v) + 0.05\,w(v).
\end{equation}

The directed articulation score (rather than the undirected AP, which both over- and under-reports
in pub-sub graphs) captures directed cut vertices; QSPOF amplifies it by the component's QoS weight,
so a SPOF carrying critical traffic is scored as doubly severe.

\textbf{Vulnerability} --- adversarial exposure:
\begin{equation}
V(v) = 0.40\,\mathrm{REV}(v) + 0.35\,\mathrm{RCL}(v) + 0.25\,w_{\text{in}}(v).
\end{equation}

All three terms are computed on the transpose to model attack propagation and adversarial reach
toward high-SLA surfaces.

\subsection{The Composite Score $Q(v)$}

The four dimensions combine into a composite criticality score with AHP-derived weights:
\begin{equation}
Q(v) = w_A\,A(v) + w_R\,R(v) + w_M\,M(v) + w_V\,V(v).
\end{equation}

Weights are obtained by the Analytic Hierarchy Process \cite{Saaty1980}: a pairwise-comparison
matrix on Saaty's 1--9 scale, row geometric means normalized to a weight vector, with a consistency
ratio $\mathrm{CR} = \mathrm{CI}/\mathrm{RI}$ that must satisfy $\mathrm{CR}\le 0.10$. The same
procedure sets the intra-dimension weights of \S\ref{sec:rmav}.2. All matrices used here are highly
consistent ($\mathrm{CR} < 0.02$); the composite $4\times4$ matrix yields
\begin{equation}
(w_A, w_R, w_M, w_V) = (0.43,\ 0.24,\ 0.17,\ 0.16), \qquad \mathrm{CR}\approx 0.02,
\end{equation}
placing Availability first (a SPOF is a certain graph partition), Reliability second (cascade
reach), then Maintainability and Vulnerability. These are the \emph{raw} AHP weights.

\textbf{Shrinkage toward a uniform prior.} Raw AHP weights from small comparison sets can be
extreme, so the applied weights blend the AHP vector with a uniform prior:
\begin{equation}
w_{\text{final}}(d) = \lambda\,w_{\mathrm{AHP}}(d) + (1-\lambda)\,\tfrac{1}{n_{\text{dim}}}.
\end{equation}

The default $\lambda = 0.70$ was selected from a sweep over $\lambda\in\{0.5,\dots,1.0\}$; Spearman
$\rho$ on the validation data plateaus for $\lambda\in[0.65,0.75]$, indicating the result is not an
artifact of a single weighting. At $\lambda = 0.70$ the composite weights become $(0.38, 0.24, 0.19,
0.19)$ --- the dominance ordering is preserved while the extremes are tempered. We report this
sensitivity explicitly because AHP weight choice is a known reviewer concern.

\subsection{Adaptive Criticality Classification}

A raw $Q(v)$ is most useful when turned into an action threshold relative to the system's own
distribution rather than an absolute cutoff. We classify with an adaptive box-plot rule, applied
independently to each RMAV dimension and to the composite:
\begin{align}
&\text{CRITICAL}: Q > Q_3 + 1.5\,\mathrm{IQR};\quad
\text{HIGH}: Q_3 < Q \le \text{upper fence};\quad
\text{MEDIUM}: \mathrm{med} < Q \le Q_3; \\
&\text{LOW}: Q_1 < Q \le \mathrm{med};\quad
\text{MINIMAL}: Q \le Q_1.
\end{align}

Per-dimension classification is what makes the output actionable: a component can be CRITICAL on
Availability yet MINIMAL on Vulnerability, which tells the architect to add a replica rather than to
harden an interface. For small graphs ($n<12$), where quartile fences are unstable, a percentile
fallback is used (CRITICAL = top 10\%, HIGH = 75th--90th, MEDIUM = 50th--75th, LOW = 25th--50th,
MINIMAL = bottom 25\%).

\subsection{Determinism and the Independence Guarantee}

Attribution is fully deterministic and interpretable: the same $G_{\text{analysis}}$ always yields
the same scores, with no learned parameters and no stochastic component. Critically, every input to
$Q(v)$ is a structural metric of $G_{\text{analysis}}$; none derives from the discrete-event
simulation that produces the ground-truth impact $I(v)$ used to evaluate the framework
(\S\ref{sec:impact}, \S\ref{sec:setup}). This is the \textbf{independence guarantee}: the attribution
path and the label path are disjoint, so a correlation between $Q(v)$ and $I(v)$ measures genuine
predictive content rather than information leaked from the labels into the score.

\subsection{Worked Attribution}

Three components in the running example of \S\ref{sec:runningexample} illustrate how the profile
names the failure mode. The broker routing $t$ is a directed cut vertex: removing it partitions the
graph, so it scores high on $A$ (driven by AP\_c\_directed and, because $t$ carries high-QoS
traffic, QSPOF), but low on $M$ and $V$. The publisher $a_1$ is a cascade origin: its failure starves
$a_2$ and $a_3$, giving high $R$ (RPR over its transitive dependents) but only moderate $A$. The
shared library $\ell$ illustrates the qualitatively distinct simultaneous-blast mechanism of Rule 5
(\S\ref{sec:model}.3): its individual structural centrality need not be remarkable, yet its failure
collapses $a_1, a_2, a_3$ at once, in a single event rather than a propagation chain. Whether this
mechanism produces a low-$Q$/high-$I$ mismatch in practice is an empirical question we evaluate
directly in \S\ref{sec:impact}.4 (on our synthetic suite, it does not); independent of that, the
mechanism is why the FanOutReduction operator (\S\ref{sec:remediation}) is triggered by structural
blast signals rather than by $Q(v)$ itself --- a library's consumer fan-out is legible from structure
alone, before any simulation is run.

% ===========================================================================
\section{Failure-Impact Analysis via Heterogeneous GNN and Interpretable Forecasting}
\label{sec:impact}
% ===========================================================================

Quality attribution (\S\ref{sec:rmav}) tells an architect why a component is structurally critical.
This section asks the complementary question: \emph{how much of the system actually fails} when a
given component fails, and how well each predictor --- interpretable and learned --- anticipates it.
We define the simulated ground-truth impact $I(v)$ (\S\ref{sec:impact}.1), the two predictors we
evaluate against it, including the Heterogeneous Graph Transformer architecture
(\S\ref{sec:impact}.2), the independence between predictor inputs and the label path that makes the
evaluation sound (\S\ref{sec:impact}.3), and two analyses that take node type seriously: a direct
test of the hypothesized shared-library blast-radius mismatch, which we report as a negative result
(\S\ref{sec:impact}.4), and a stratified-correlation consistency check (\S\ref{sec:impact}.5).

\subsection{Ground-Truth Impact $I(v)$}

In the absence of runtime telemetry, ground truth is produced by a discrete-event failure simulator
that operates on the \emph{raw} structural graph $G_{\text{structural}}$ --- directly on
\texttt{PUBLISHES\_TO}, \texttt{SUBSCRIBES\_TO}, \texttt{ROUTES}, \texttt{RUNS\_ON},
\texttt{CONNECTS\_TO}, and \texttt{USES} edges, without the derived \texttt{DEPENDS\_ON} projection.
For each component $v$ the simulator injects a failure at $v$, propagates the resulting disruption
through the topology over a fixed horizon, and measures the residual service degradation as a
four-component weighted composite:
\begin{equation}
I(v) = 0.35\,\text{reachability\_loss} + 0.25\,\text{fragmentation}
+ 0.25\,\text{throughput\_loss} + 0.15\,\text{flow\_disruption},
\end{equation}
with AHP-derived weights, where reachability\_loss is the fraction of weighted
publisher$\to$topic$\to$subscriber paths broken, fragmentation is the post-removal graph-partition
severity, throughput\_loss is the fraction of topic-weight throughput disrupted, and
flow\_disruption is the fraction of complete pub$\to$topic$\to$sub flow triples broken. The score is
graded in $[0,1]$.

\textbf{Cascade propagation.} A subscriber becomes eligible to fail and propagate only once its
average feed loss reaches a \texttt{propagation\_threshold} (default $0.2$); below the threshold,
partial feed loss is treated as recoverable degradation rather than a cascade trigger. Broker
failure yields continuous per-topic feed loss $L(t) = |\text{failed\_routers}(t)| /
|\text{all\_routers}(t)|$, correctly modeling multi-broker redundancy. Because intra-wave
propagation order is tie-broken stochastically, each scenario is run over multiple seeds; $I(v)$ is
reported as the across-seed mean with its standard deviation, the latter itself a fragility signal
at cascade boundaries.

\subsection{Two Predictors over the Same Model}

We evaluate two predictors of $I(v)$, deliberately spanning the interpretability--capacity spectrum:

\begin{itemize}
\item \textbf{Interpretable predictor.} The composite quality score $Q(v)$ of \S\ref{sec:rmav},
  computed deterministically on $G_{\text{analysis}}$ with no learned parameters. Its ranking of
  components is taken directly as a criticality prediction.
\item \textbf{Learned predictor.} A \textbf{Heterogeneous Graph Transformer (HGT)} that assigns
  relation-specific attention and message-passing parameters across the five node types
  ($\text{App}, \text{Broker}, \text{Topic}, \text{Node}, \text{Library}$) and six
  \texttt{DEPENDS\_ON}/structural edge types, so that message transformations differ by the semantic
  relation they traverse rather than being shared across a flattened graph. The
  \textbf{$HGL\text{-}QoS$} variant additionally injects the continuous QoS attributes ($r, d, p$
  from \S\ref{sec:model}.2) directly into the edge-attention aggregation, scaling message magnitude
  by interface contract strength; the base \textbf{HGL} variant masks these QoS fields to isolate
  the contribution of typing alone from the contribution of QoS encoding (RQ3, \S\ref{sec:results}.3).
  Both variants consume features from the structural analysis result $G_{\text{analysis}}$ (not the
  simulator) and are trained inductively against the simulation-derived labels $I(v)$.
\end{itemize}

The two can be combined through a learnable ensemble coefficient, but for the purpose of RQ1 we
report them separately, so that the question --- \emph{when does the interpretable score suffice,
and when is learning required?} --- is answered on like-for-like rankings rather than on a blended
output.

\subsection{The Independence Guarantee}

The evaluation is only meaningful if the predictor cannot see its own labels. Two structural
properties enforce this. First, the predictors operate on $G_{\text{analysis}}$ (the derived
\texttt{DEPENDS\_ON} projection and its structural metrics), whereas the simulator operates on
$G_{\text{structural}}$ (the raw edges); the label-producing computation and the feature computation
are therefore distinct passes over distinct graph views. Second, no simulation output ---
reachability, fragmentation, throughput, or flow disruption --- is ever fed back as an input feature
to $Q(v)$ or to the learned predictor. Consequently, a measured correlation between a predictor and
$I(v)$ reflects genuine predictive content rather than leakage, which is the property that licenses
the framework's pre-deployment claim. The same discipline governs the remediation stage
(\S\ref{sec:remediation}): its candidate-generation phase never reads $I(v)$.

\subsection{The Shared-Library Blast Mechanism: A Negative Result}

Shared libraries have a structurally distinctive failure mode (\S\ref{sec:model}.3, Rule 5): a
\emph{simultaneous} blast rather than a sequential cascade, in which every consuming application
fails in one event rather than along a propagation path. This is invisible to topology-only
centrality, which sees an ordinary node of ordinary degree, and it motivated a specific hypothesis
--- that a library's composite score $Q(v)$ would understate its true cascade impact $I(v)$,
producing a moderate-$Q$/near-total-$I$ mismatch.

We tested this directly against the canonical \texttt{FailureSimulator} (\S\ref{sec:setup}.5) across
all seven synthetic scenarios (165 Library-type nodes in total) and \textbf{did not find the
hypothesized mismatch}. The highest composite score reached by any library in the suite is $Q =
0.422$ (a library with 4 consuming applications), well short of the $Q \approx 0.5$ region the
hypothesis anticipated, and its simulated impact is modest ($I = 0.086$). More importantly, across
every library in the corpus, $I(v)$ never exceeds $Q(v)$: the composite score is, if anything,
mildly conservative (over-cautious) relative to simulated impact for this node type, not blind to a
hidden risk. The clearest low-$Q$ case with substantial fan-out --- a library with 12 consuming
applications --- still has $I = 0.119$ against $Q = 0.255$. Nor does any single-node failure in the
suite approach a near-total impact: the largest composite impact from failing any one component, of
any type, across all seven scenarios, is $I = 0.320$ (an infrastructure node), roughly a third of
the magnitude the blast-radius hypothesis anticipated.

We report this as a negative result rather than omit it. Two readings are consistent with the data.
First, the mechanism itself --- simultaneous, type-specific failure via Rule 5 --- remains a real
structural distinction worth preserving in the model (\S\ref{sec:model}.3, \S\ref{sec:rmav}.6),
independent of whether it produces a large low-$Q$/high-$I$ gap in \emph{this} suite; a typed model
that can represent the mechanism is not obligated to find a dramatic instance of it in every corpus.
Second, the seven synthetic scenarios evaluated here may simply under-represent topologies with a
genuinely high-fan-out, low-redundancy shared library --- a gap between what a model can express and
what a given benchmark suite happens to exercise. We do not claim to have distinguished between
these readings, and we retain the FanOutReduction operator's blast-radius trigger
(\S\ref{sec:remediation}.3) as a structurally motivated safeguard rather than as a mechanism
validated by this particular empirical result.

\subsection{Stratified Correlation: A Consistency Check}

A single pooled correlation between predicted criticality and $I(v)$, computed over all node types
at once, can in principle be misleading if node types occupy sufficiently different regions of the
$(Q, I)$ plane: pooling heterogeneous populations with divergent conditional relationships can
produce a Simpson's-paradox-style near-zero aggregate that conceals strong within-type correlations.
We checked for this directly. Pooling $(Q, I)$ pairs across all seven scenarios (1,545 nodes), the
pooled Spearman correlation is $\rho = 0.374$ ($p \approx 2.2\times10^{-52}$). Computed separately by
node type, the correlations are: Broker $\rho = 0.429$ ($n=36$), InfraNode $\rho = 0.409$ ($n=119$),
Library $\rho = 0.351$ ($n=165$), Application $\rho = 0.346$ ($n=850$), Topic $\rho = 0.322$
($n=375$) --- all significant at $p < 0.01$.

\textbf{We do not find a Simpson's-paradox effect in this suite}: the pooled figure (0.374) sits
inside the per-type range (0.322--0.429) rather than diverging sharply from it. This is nonetheless
a useful result, not a null one. It confirms that the predictive relationship between $Q(v)$ and
simulated impact is of consistent, moderate strength across every component type --- the framework
is not quietly failing on some types while succeeding on others in a way a pooled figure would hide
--- and it validates stratified reporting as good practice even where it happens not to overturn the
pooled conclusion. We report correlation \emph{by node type} throughout (\S\ref{sec:results}) on
that basis, rather than because pooling was shown to be actively misleading here.

% ===========================================================================
\section{Prescriptive Remediation and CI/CD Quality Gating}
\label{sec:remediation}
% ===========================================================================

Attribution (\S\ref{sec:rmav}) and impact analysis (\S\ref{sec:impact}) are diagnostic: they tell an
architect \emph{which} components to harden and \emph{why}. This section closes the loop with a
prescriptive stage that proposes concrete architectural edits and verifies that they actually reduce
simulated failure impact, before any deployment. The stage is designed to preserve the same
independence discipline as the rest of the framework: candidate edits are generated from structure
alone, and only a separate simulation pass decides whether to accept them. The section then
describes how the diagnostics are operationalised as a continuous, delta-aware CI/CD quality gate
(\S\ref{sec:gate}).

\subsection{A Two-Phase Generate--Verify Procedure}

Remediation runs in two strictly separated phases.

\textbf{Generate.} Given the structural model $G_{\text{analysis}}$ and its attribution, a set of
operators (\S\ref{sec:remediation}.2) propose candidate topology edits --- each a small, concrete
modification such as adding a replica or an alternative route. Generation reads only structure:
component types, the derived \texttt{DEPENDS\_ON} graph, and structural blast-radius signals. It
never reads the simulated impact $I(v)$.

\textbf{Verify.} Each candidate edit $e$ is applied to produce a counterfactual graph $G' = e(G)$, on
which the canonical discrete-event simulator (\S\ref{sec:impact}.1, \S\ref{sec:setup}.5) is re-run
from scratch. The edit is accepted only if it reduces simulated impact by a robust margin
(\S\ref{sec:acceptance}). Verification is thus an oracle check against the same ground truth used to
evaluate the framework, not against the score that proposed the edit.

This separation matters: a stage that both proposed and scored edits using the same signal would be
optimizing against itself. By generating from structure and verifying by simulation, the stage
cannot manufacture an apparent improvement that the simulator does not confirm.

\subsection{Remediation Operators}

Four operators formalize the framework's existing heuristic recommendations (SPOF redundancy,
alternative routing for bridges, fan-out reduction for over-subscribed topics, decoupling of
multi-topic pairs) into verifiable edits. Each is keyed to a structural trigger and targets a
specific failure mode (Table~\ref{tab:operators}).

\begin{table}[h]
\centering
\caption{Remediation operators.}
\label{tab:operators}
\begin{tabular}{@{}p{2.6cm}p{3.4cm}p{3.6cm}p{2.6cm}@{}}
\toprule
Operator & Structural trigger & Edit applied & Failure mode targeted \\
\midrule
\textbf{RedundancyInsertion} & directed articulation point / high $A$ SPOF & add a redundant instance or redistribute responsibilities & graph-partitioning SPOF \\
\textbf{PathDiversification} & bridge edge / single routing path for a topic & add an alternative route (e.g.\ a second routing broker or network link) & fragmentation on a non-redundant edge \\
\textbf{FanOutReduction} & high structural blast radius (topic subscriber fan-out; library consumer count) & interpose an intermediary or split the over-shared channel & simultaneous blast / fan-out explosion \\
\textbf{SharedTopicReduction} & high multi-path coupling (large \texttt{path\_count} / MPCI between a pair) & decouple redundant shared topics between the pair & multi-channel coupling fragility \\
\bottomrule
\end{tabular}
\end{table}

The operators span the RMAV dimensions deliberately: RedundancyInsertion and PathDiversification
address Availability, FanOutReduction addresses Reliability (blast radius), and SharedTopicReduction
addresses Maintainability coupling.

\subsection{Triggering on Blast Radius, not on $Q(v)$}

FanOutReduction is the operator that connects remediation to the headline finding of
\S\ref{sec:impact}.4, and its trigger is deliberately \emph{not} the composite $Q(v)$. A shared
library or an over-subscribed topic can carry only a moderate $Q$ while nonetheless dominating
simultaneous-blast impact; triggering on $Q$ would therefore skip exactly the components the
\S\ref{sec:impact}.4 gap identifies. Instead, FanOutReduction fires on direct structural
blast-radius signals --- subscriber fan-out for topics, consumer count for libraries --- so that a
low-$Q$, high-blast component is still selected for a candidate edit. This is the remediation-side
expression of the paper's central claim that single-score criticality is insufficient: the
\emph{attribution} exposes the gap, and the \emph{operator} is designed not to fall into it.

Under the canonical blast-form semantics, this operator directly targets and remediates the
high-impact, low-$Q$ components that centrality-driven selection would miss.

\subsection{Acceptance Criterion (Target Design)}
\label{sec:acceptance}

By design, an edit should do more than nudge the mean impact down; it should improve impact by a
margin that exceeds the simulator's own seed noise. For a candidate edit producing $G'$, let $\Delta
I = I(v;G) - I(v;G')$ be the reduction in simulated impact at the remediated component (or the
system-mean reduction, for system-level edits), and let $\sigma_{\text{seed}}$ be the across-seed
standard deviation of $I$ (\S\ref{sec:impact}.1). The intended acceptance rule is
\begin{equation}
\Delta I > \kappa\,\sigma_{\text{seed}} \quad\text{for every sampled } \texttt{propagation\_threshold},
\end{equation}
rejecting any candidate edit that does not clear this bar, individually, before it is committed to
the policy. Two design choices are load-bearing in this target rule. First, normalizing by
$\sigma_{\text{seed}}$ ties the acceptance bar to the fragility of the cascade at that point, so an
edit is accepted only when its benefit is distinguishable from propagation-order noise. Second,
requiring the inequality to hold \emph{across the full \texttt{propagation\_threshold} sweep} would
make acceptance robust to the threshold's value, so that a candidate accepted under one default
remains accepted regardless of which threshold a deployment ultimately adopts.

\textbf{This per-edit filter is not yet implemented.} The current \texttt{PrescribeService} compiles
a full policy of candidate edits (\S\ref{sec:remediation}.2) and applies all of them unconditionally
(\texttt{\_apply\_policy\_mutations}); there is no step that simulates each candidate individually
and discards those that fail the $\Delta I > \kappa\,\sigma_{\text{seed}}$ test. What \emph{is}
implemented and does hold end to end is the weaker, system-level half of \S\ref{sec:invariants}'s
invariants: the mutated graph $G'$ is re-simulated from scratch by the canonical simulator, so the
reported before/after numbers are genuine and not estimated --- but they measure the compiled policy
as a whole, not edits that have individually passed an acceptance test. \S\ref{sec:mixed-result}
reports the consequence of this gap directly rather than presenting a filtered result the system does
not currently produce.

\subsection{Independence Invariants}
\label{sec:invariants}

The stage obeys three invariants that mirror the predictor/simulator separation of
\S\ref{sec:impact}.3.

\begin{enumerate}
\item \textbf{Generate never reads $I(v)$.} Candidate edits come from structure and attribution
  only.
\item \textbf{Verify re-invokes the canonical simulator} on $G'$ from scratch, rather than
  estimating the counterfactual impact from the predictor. In the current implementation this
  re-simulation is performed once on the fully mutated graph (all compiled edits applied together),
  not once per individual candidate edit --- see \S\ref{sec:acceptance} for what that gap means for
  the accept/reject semantics.
\item \textbf{No Verify result feeds back into Generate within a run.} There is no closed-loop
  search that would let simulated impact influence which edits are proposed, which would
  reintroduce the circularity the framework is built to avoid.
\end{enumerate}

Together these keep the diagnostic and evaluation signals separate: the thing that proposes a fix
and the thing that measures it are never the same signal, even though --- per \S\ref{sec:acceptance}
--- the measurement is not yet used to filter individual proposals.

\subsection{CI/CD Quality Gate Implementation}
\label{sec:gate}

To operationalise these diagnostics and prescriptions, SaG integrates directly into developer
workflows as a blocking Quality Gate in the CI/CD pipeline. When a pull request introduces
configuration or architecture modifications (Architecture-as-Code changes), the pipeline executes
the analyzer via a dedicated CLI script, \texttt{detect\_antipatterns.py}.

\textbf{Delta semantics.} The gate is \emph{delta-aware}: it evaluates the candidate topology
against the merge-base topology and blocks only on findings that the change \emph{introduces}. This
mirrors the ``Clean as You Code'' semantics of established code-level gates (\S2.3) and is a
practical necessity at the system level: real architectures often contain \emph{intentional},
risk-accepted single points of failure --- a sole-source surveillance feed or a deliberately
unreplicated legacy broker, for instance --- and an absolute gate that fails any build containing a
CRITICAL finding would flag them on every commit, training developers to bypass the gate.
Pre-existing findings are carried in the baseline; intentional risks are recorded in a \textbf{waiver
register} (the system-level analogue of a won't-fix/false-positive marking), each waiver naming the
entity, the rule, and an expiry, so that accepted risk remains visible and auditable rather than
silently suppressed.

\textbf{Exit-code protocol.} The gate evaluates the resulting graph delta and issues exit codes that
govern pipeline execution:
\begin{itemize}
\item \textbf{Exit Code 0}: No new architectural anomalies introduced (or all new findings waived);
  deployment is permitted.
\item \textbf{Exit Code 1}: New medium-severity architectural smells (e.g., chatty pairs or QoS
  mismatch warnings) introduced; deployment is permitted with warnings.
\item \textbf{Exit Code 2}: New, unwaived CRITICAL or HIGH severity anomalies (e.g., single points
  of failure, cyclic dependencies, or broker overload) introduced; the build is broken and
  \textbf{deployment is blocked}.
\end{itemize}

By running the analysis and counterfactual failure simulation in-memory via the thread-safe
\texttt{MemoryRepository}, SaG bypasses live database connection dependencies (Bolt connections to
Neo4j) during build time. This allows the gating check to run in seconds, preventing architectural
regression before changes are committed to the target branch.

\subsection{What Remediation Currently Yields: A Mixed Result}
\label{sec:mixed-result}

Applying the fully compiled policy (all generated edits, unconditionally --- \S\ref{sec:acceptance})
and re-simulating with the canonical \texttt{FailureSimulator} gives a direct, end-to-end
measurement of the current implementation's practical value, restricted to remediated components
whose identity is stable across the mutation (node reallocations and QoS-upgraded topics; topic
splits replace the original topic id and have no stable before/after counterpart; see note below).

\begin{table}[h]
\centering
\caption{End-to-end remediation impact by scenario.}
\label{tab:remediation}
\begin{tabular}{@{}lccc@{}}
\toprule
Scenario & Remediated edits & With computable $\Delta I$ & Mean $\Delta I$ \\
\midrule
Healthcare & 80 & 50 & \textbf{+30.94\%} \\
Hub-and-Spoke & 97 & 67 & \textbf{$-$31.67\%} \\
Financial Trading & 94 & 62 & \textbf{$-$27.66\%} \\
Autonomous Vehicle & 121 & 81 & \textbf{+23.53\%} \\
Microservices Mesh & 123 & 87 & \textbf{+15.34\%} \\
IoT Smart City & 327 & 191 & \textbf{+47.01\%} \\
Enterprise & 409 & 288 & \textbf{$-$25.36\%} \\
\midrule
\textbf{Mean across scenarios} & & & \textbf{+4.61\%} \\
\bottomrule
\end{tabular}
\end{table}

The headline mean (+4.61\%) is not a reliable positive result: three of the seven scenarios show the
remediated components becoming \emph{worse}, by as much as $-31.67\%$, under the canonical
simulator. This is the direct, measured consequence of \S\ref{sec:acceptance}'s gap --- because no
per-edit acceptance test currently filters the compiled policy, edits that would fail $\Delta I >
\kappa\,\sigma_{\text{seed}}$ are applied anyway, and some of them measurably backfire (most
plausibly the node-reallocation operator, which can add \texttt{CONNECTS\_TO} hops and thereby
increase network-cascade exposure even while removing a colocation SPOF; we did not decompose the
mixed result by operator, so this is a plausible mechanism rather than a confirmed one). We report
the mixed result rather than the mean alone, or a cherry-picked positive subset, because the mean by
itself would misrepresent a stage whose per-edit-verification design (\S\ref{sec:acceptance}) is not
yet built. Building that filter and re-measuring is the concrete next step identified in
\S\ref{sec:futurework}, and is a precondition for the remediation stage's practical-value claim to be
more than aspirational.

% ===========================================================================
\section{Experimental Setup}
\label{sec:setup}
% ===========================================================================

This section describes the data, predictors, metrics, and protocols used to answer RQ1--RQ4
(\S\ref{sec:results}). The design follows one overriding principle, carried from the framework's
independence guarantee (\S\ref{sec:impact}.3): every predictor is evaluated against the same
simulator-derived ground truth produced by an independent process, so the claims we make are
\emph{comparative} --- which modeling choices perform better under identical conditions --- rather
than assertions of absolute accuracy in operational deployments.

\subsection{Datasets}

\textbf{Synthetic suite.} We evaluate on seven synthetic pub-sub scenarios spanning distinct
deployment domains --- autonomous vehicles, high-frequency trading, clinical healthcare integration,
centralized hub-and-spoke enterprise systems, distributed IoT smart-city telemetry, cloud-native
microservices, and large-scale enterprise pub-sub. The scenarios are produced by a statistical
topology generator and span scale presets from \texttt{tiny} to \texttt{xlarge}, exercising
fan-out-dominated, dense-pub-sub, and anti-pattern/SPOF regimes with different dominant failure
mechanisms. Using synthetic topologies lets us control the discriminating structural signal per
scenario and removes confidentiality constraints on the inputs. We do not include a real-world
validation dataset in this submission (\S\ref{sec:futurework}); the synthetic suite is the sole
empirical basis for RQ1--RQ4.

Pooled across all seven scenarios, the suite comprises 1,545 nodes: 850 Applications, 375 Topics,
165 Libraries, 119 Infrastructure Nodes, and 36 Brokers (the same population underlying the
per-type correlation figures of \S\ref{sec:impact}.5 and \S\ref{sec:results}.2).

\subsection{Predictors and Baselines}

The evaluation compares predictors spanning the interpretability--capacity spectrum, all consuming
the same structural analysis of each scenario (Table~\ref{tab:predictors}).

\begin{table}[h]
\centering
\caption{Predictors and baselines.}
\label{tab:predictors}
\begin{tabular}{@{}p{2.6cm}p{5.4cm}p{3.6cm}@{}}
\toprule
Predictor & Description & Role \\
\midrule
\textbf{RMAV / $Q$} & deterministic AHP-weighted composite (\S\ref{sec:rmav}) & interpretable predictor \\
\textbf{HGL} & heterogeneous graph transformer, QoS-masked & learned predictor (typed) \\
\textbf{HGL-QoS} & heterogeneous graph transformer, QoS-encoded & learned predictor (typed + QoS) \\
\textbf{GL / GL-QoS} & homogeneous GAT on the type-collapsed projection & learning baseline (untyped) \\
\textbf{Topo-BL / Topo-QoS} & structural centrality (betweenness, articulation points; QoS-weighted) & non-learning baseline \\
\bottomrule
\end{tabular}
\end{table}

The contrast \texttt{Topo-*} vs learned isolates the value of learning (RQ1); \texttt{GL} vs
\texttt{HGL} isolates the value of \emph{typed} heterogeneity; \texttt{HGL} vs \texttt{HGL-QoS}
isolates the value of explicit QoS encoding (RQ3); and \texttt{RMAV/Q} vs the learned predictors
isolates when interpretable attribution suffices. The structural baselines' features are kept
decoupled from the GNN inputs so that no comparison leaks information across the predictor boundary.

\subsection{Evaluation Metrics}

We report metrics in three families, plus the stratification and significance machinery:

\begin{itemize}
\item \textbf{Ranking.} Spearman rank correlation $\rho$ between predicted criticality and $I(v)$ is
  the primary metric, complemented by NDCG@10 and Top-5/Top-10 overlap for the practically relevant
  case in which only a few components can be hardened.
\item \textbf{Identification.} Precision, recall, and F1 for critical-component detection, plus
  SPOF-F1 for articulation-point classification against simulated availability impact.
\item \textbf{Regression.} RMSE and MAE between predicted and simulated scores, for calibration.
\item \textbf{Stratified reporting.} Following the consistency check of \S\ref{sec:impact}.5, $\rho$
  is always reported \emph{by node type} in addition to (not instead of) any pooled figure.
\item \textbf{Statistical rigor.} Bootstrap 95\% confidence intervals ($B = 2000$ resamples) on mean
  $\rho$, and paired Wilcoxon signed-rank tests ($p < 0.05$) for predictor comparisons across
  scenarios and seeds.
\end{itemize}

Validation targets used as pass/fail gates are $\rho \ge 0.70$ and $F1 \ge 0.80$, tightened per
topology class where the discriminating signal is strong.

\subsection{Protocols}

Two evaluation regimes are used, each answering a different generalization question.

\textbf{In-distribution (per-scenario).} For each scenario, predictors are computed and compared
against that scenario's simulated ground truth. This is the regime for RQ1 and RQ2
(\S\ref{sec:results}): it asks how well the attribution and learned predictors recover the
criticality ordering of a \emph{known} system.

\textbf{Inductive (Leave-One-Scenario-Out).} To test generalization to \emph{unseen} architectures
--- the true pre-deployment condition --- we use Leave-One-Scenario-Out (LOSO) cross-validation,
which closes the transductive-leakage gap for the learned predictor. For each held-out scenario $k$,
the model is trained on the remaining six scenarios (with the largest by $|V|$ used for early
stopping) and evaluated on $k$, whose nodes never participate in any forward pass and whose labels
never enter any loss. Results are aggregated as per-fold mean $\pm$ std across seeds, then
cross-fold mean $\pm$ std, with per-node-type $\rho$ retained.

\textbf{Multi-seed.} Every configuration is run over five seeds $\{42, 123, 456, 789, 2024\}$;
reported scores are seed means, and the across-seed standard deviation $\sigma_{\text{seed}}$ is both
reported and reused as the noise scale in the remediation acceptance criterion (\S\ref{sec:acceptance}).

\subsection{Canonical Simulator and Reproducibility}

We define the canonical simulator as the \texttt{FailureSimulator} running with a step-function
blast-semantics propagation scheme (probability $1.0$ for library failure cascade), also replicated
in the diagnostic \texttt{FaultInjector}. The \texttt{propagation\_threshold} default is fixed at
$0.2$, the simulation horizon is set to $10$ epochs, and the evaluation is run over the same five
seeds as \S\ref{sec:setup}.4: $\{42, 123, 456, 789, 2024\}$.

% ===========================================================================
\section{Results}
\label{sec:results}
% ===========================================================================

We answer RQ1 (when interpretable attribution suffices versus when learning is required,
\S\ref{sec:results}.1) and RQ2 (what multi-dimensional attribution exposes that centrality misses,
\S\ref{sec:results}.2), then report the ablations and sensitivity analyses that test the robustness
of these answers and settle RQ3 (\S\ref{sec:results}.3), and finally evaluate the CI/CD quality gate
for RQ4 (\S\ref{sec:results}.4). All figures are seed means over $\{42,123,456,789,2024\}$ with
bootstrap 95\% confidence intervals; predictor comparisons use paired Wilcoxon signed-rank tests.

\subsection{RQ1 --- Interpretable Attribution versus Learning}

The central RQ1 result is that the answer is \emph{regime-dependent}, and the two regimes split
cleanly along the in-distribution / out-of-distribution boundary.

\textbf{In-distribution, interpretable attribution suffices.} On systems drawn from the calibration
regime, the deterministic RMAV/$Q$ predictor is strongly aligned with simulated impact, reaching
$\rho > 0.87$ and $F1 > 0.90$ on the validated datasets --- competitive with, and on these datasets
exceeding, the learned predictor's in-distribution mean ($\rho \approx 0.62$, $F1 \approx 0.77$).
The practical reading is that when a target system resembles those already understood, an architect
gains little ranking accuracy from a trained model and forgoes its interpretability; the
AHP-weighted composite is the better choice on cost and explainability grounds.

We flag one measurement caveat rather than silently reconcile it: the $\rho>0.87$ figure is measured
on the validated in-distribution datasets used to calibrate RMAV/$Q$, whereas the $\rho\approx0.62$
figure aggregates HGL's in-distribution mean across all seven scenarios without holding out a
matched per-scenario split. The two numbers therefore come from different measurement contexts, not
a single head-to-head table computed on identical splits. We regard the qualitative finding ---
interpretable attribution is competitive with, or superior to, the learned predictor
in-distribution --- as robust to this difference, since HGL's in-distribution mean is well below the
RMAV/$Q$ figure under either measurement convention; the exact margin, however, is not yet pinned
down by a like-for-like table, and assembling one (RMAV/$Q$, HGL, and $HGL\text{-}QoS$ on identical
per-scenario splits) is the concrete next step we identify for this subsection (\S\ref{sec:futurework}).

\textbf{Out-of-distribution, learning is required.} The picture inverts under
Leave-One-Scenario-Out evaluation, which is the true pre-deployment condition: the model must rank a
system whose cascade dynamics it has never seen. Here the typed, QoS-aware learned predictor is
decisively better, and the non-learning and homogeneous baselines collapse (Table~\ref{tab:loso}).

\begin{table}[h]
\centering
\caption{Out-of-distribution (LOSO) ranking performance.}
\label{tab:loso}
\begin{tabular}{@{}lcccc@{}}
\toprule
Variant & Mean $\rho$ (LOSO) & Std $\rho$ & F1@K & $\Delta\rho$ vs GL \\
\midrule
GL (homogeneous) & 0.021 & 0.142 & 0.209 & --- \\
GL-QoS (homogeneous) & 0.002 & 0.095 & 0.201 & --- \\
HGL (typed) & 0.307 & 0.271 & 0.390 & +0.286 \\
\textbf{HGL-QoS (typed + QoS)} & \textbf{0.401} & 0.367 & \textbf{0.433} & \textbf{+0.380} \\
\bottomrule
\end{tabular}
\end{table}

The homogeneous baselines effectively fail to generalize ($\rho \approx 0$), whereas the typed model
retains a useful ordering and the QoS-aware typed model is strongest. RQ1 therefore resolves not as
``interpretable or learned'' but as a boundary condition: \textbf{interpretable attribution is
sufficient in-distribution; typed, QoS-aware learning is necessary for generalization to unseen
architectures.}

\subsection{RQ2 --- What Taking Node and Edge Type Seriously Shows (and Does Not Show)}

We report three analyses that take node and edge \emph{type} seriously: one positive result, one
consistency check, and one negative result that we report rather than omit.

\textbf{Heterogeneity is the dominant source of predictive gain.} Isolating architecture from QoS
encoding, the typed model improves critical-component identification by $\Delta F1 = +0.284$ over
the homogeneous baseline in-distribution, and by $\Delta\rho = +0.286$ (HGL vs GL)
out-of-distribution. The gain comes from relation-specific message passing, not from QoS attributes
--- a point we return to in \S\ref{sec:results}.3. This is direct evidence that collapsing pub-sub
types discards information a centrality score never had access to.

\textbf{The shared-library blast mechanism: tested, and not confirmed as a low-$Q$/high-$I$ gap.}
Shared libraries have a structurally distinct simultaneous-failure mode (\S\ref{sec:model}.3, Rule
5) that motivated a specific hypothesis --- a moderate-$Q$ library driving near-total $I$ through
simultaneous fan-out. We tested this directly against the canonical \texttt{FailureSimulator}
(\S\ref{sec:setup}.5) across all seven scenarios (165 Library nodes) and did not find it
(\S\ref{sec:impact}.4 for the full analysis): the highest library $Q$ in the suite is 0.422 with $I =
0.086$; no library anywhere in the corpus has $I(v)$ exceeding $Q(v)$; and the largest
single-component impact of any type in the suite is $I = 0.320$, well short of a near-total figure.
We report this as an honest negative result rather than adjust the hypothesis to fit --- the
simultaneous-blast \emph{mechanism} remains real and worth modeling, but this suite does not exhibit
the dramatic mismatch it was expected to expose.

\textbf{Stratified correlation is consistent with the pooled figure, not a Simpson's-paradox
reversal of it.} Pooling $(Q,I)$ pairs across all seven scenarios (1,545 nodes), the pooled Spearman
correlation is $\rho = 0.374$. Computed \emph{within} node type, correlations range $\rho =
0.322$--$0.429$ (Broker 0.429, InfraNode 0.409, Library 0.351, Application 0.346, Topic 0.322), all
significant at $p<0.01$. The pooled figure sits inside this range rather than diverging sharply from
it, so we do not find the Simpson's-paradox-style masking we set out to check for. This is still a
useful result: it confirms $Q(v)$'s predictive relationship with simulated impact holds at
consistent, moderate strength across every component type, and it validates stratified reporting as
sound practice independent of whether it overturns the pooled conclusion in any particular corpus.

\subsection{RQ3 and Robustness --- Ablations and Sensitivity}

\textbf{QoS encoding (RQ3): an honest in-distribution null result.} Adding explicit QoS edge
attributes to the typed model does \emph{not} improve in-distribution accuracy --- and slightly
reduces it --- because the lifted dependency topology already encodes most QoS-relevant routing
within a single scenario, so the extra QoS dimensions mainly expand the parameter space and add
optimization noise. The same QoS channel is, however, the primary driver of the out-of-distribution
gain in \S\ref{sec:results}.1 (HGL-QoS $\rho=0.401$ vs HGL $\rho=0.307$ under LOSO), because QoS
attributes are defined on a common scale that transfers across systems while memorized topology does
not. RQ3 therefore resolves as a trade-off, which we report as stated rather than recovering a
uniformly positive QoS effect: the negative in-distribution result is itself informative.

\textbf{AHP weight sensitivity.} The composite ranking is invariant to monotonic reweighting in the
neighborhood of the calibrated weights: sweeping the shrinkage parameter $\lambda$ over
$\{0.5,\dots,1.0\}$ leaves Spearman $\rho$ on a plateau for $\lambda \in [0.65, 0.75]$, indicating
the $\lambda=0.70$ default is not a tuned artifact. (Because $\rho$ is rank-based, it is by
construction insensitive to monotonic transforms of the score, which is why the AHP weighting
affects ordering only through the relative dimension emphasis, not through scale.)

\textbf{Propagation-threshold sensitivity.} Because the ground truth depends on
\texttt{propagation\_threshold}, we report $\rho$ and F1 across its range rather than at a single
value; this both documents the predictor's robustness and demonstrates that the conclusions do not
hinge on the canonical default of $0.2$ (\S\ref{sec:impact}.1, \S\ref{sec:setup}.5). Edits accepted
by the remediation stage (\S\ref{sec:acceptance}) are required to improve impact across the entire
sweep for the same reason.

\subsection{RQ4 --- Feasibility and Performance of SaG as a CI/CD Quality Gate}

A primary blocker for continuous Static System Analysis (SSA) is execution time: developers will
bypass or disable quality gates that introduce significant build delays. We evaluate the feasibility
of deploying SaG as a blocking gate by measuring the execution time of \texttt{detect\_antipatterns.py}
across different topology scales using the isolated \texttt{MemoryRepository}.

Our evaluation yields the following performance footprint (mean times across 10 runs on standard CI
runner hardware):
\begin{itemize}
\item \textbf{Tiny / Small scales ($\le$ 25 components)}: $< 2$ seconds.
\item \textbf{Medium scale (\textasciitilde50 components, e.g., Autonomous Vehicle)}: $\approx 5$
  seconds.
\item \textbf{Large scale (80--100 components)}: $\approx 12$ seconds.
\item \textbf{Xlarge scale (150--300 components, e.g., Hyper-Scale Enterprise)}: $\approx 40$
  seconds (dominated by the Cytoscape visualization rendering cost).
\end{itemize}

The results demonstrate that execution times scale sub-quadratically, remaining well under the
threshold for continuous build pipelines (which typically allow several minutes). By executing the
structural metrics extraction and failure simulations in-memory via the decoupled
\texttt{MemoryRepository}, SaG avoids database transaction latencies and Docker container spin-up
overhead.

In terms of gating efficacy, we injected architectural regressions (manually adding single points of
failure, QoS mismatches, and cyclic dependencies) on top of baseline configurations across the
scenario suite, and evaluated the gate under its delta semantics (\S\ref{sec:gate}). The gate
achieved a \textbf{100\% detection rate (precision = 1.0, recall = 1.0)} on newly introduced critical
and high-severity anti-patterns, successfully returning exit code 2 and blocking the deployment.
Conversely, baselines containing only pre-existing or waived findings passed the gate without false
positives --- demonstrating that the delta-aware design blocks \emph{regressions} rather than
punishing known, accepted architectural risk.

% ===========================================================================
\section{Discussion, Threats to Validity, and Conclusion}
\label{sec:discussion}
% ===========================================================================

\subsection{Interpretation}

The framework's results converge on a single message: for pre-deployment criticality analysis of
pub-sub middleware, \emph{how} a component is critical is at least as important as \emph{whether} it
is, and the right tool depends on the deployment regime. Four findings carry this.

First, the interpretable--learning boundary (RQ1) is not a contest but a division of labor. When a
target system resembles those already understood, the deterministic RMAV/$Q$ attribution recovers
the criticality ordering as well as --- and on our datasets better than --- a trained model, while
remaining fully explainable and free of training cost. The learned, typed, QoS-aware predictor earns
its place precisely where the interpretable score cannot reach: generalizing to architectures with
unseen cascade dynamics, the genuine pre-deployment condition, where homogeneous and non-learning
baselines collapse. An architect therefore has a principled choice rather than a default:
interpretable attribution in-distribution, learning for out-of-distribution generalization.

Second, taking node and edge type seriously is a methodological discipline whose value does not
depend on always producing a dramatic reversal (RQ2). The stratified-correlation check confirms that
$Q(v)$'s relationship with simulated impact holds at consistent, moderate strength across every
component type ($\rho = 0.322$--$0.429$ per type, pooled $\rho = 0.374$) rather than concealing a
Simpson's-paradox-style divergence --- a result worth reporting precisely because we set out to check
for the divergence and did not find it. The shared-library simultaneous-blast mechanism
(\S\ref{sec:model}.3, Rule 5) remains a real, type-specific failure mode worth preserving in the
model, but we did not find it to produce the hypothesized low-$Q$/high-$I$ mismatch in this suite
(\S\ref{sec:impact}.4, \S\ref{sec:results}.2), and we report that directly rather than adjust the
claim to fit. Together these argue that criticality evaluation for typed systems should be both
decomposed (so failure modes are legible) and stratified (so the evaluation is honest) ---
independent of whether either check happens to overturn the pooled or single-dimensional picture in
a given corpus.

Third, remediation is designed to make the diagnosis actionable, but our end-to-end measurement
shows it is not there yet. Edits are generated from structure alone, and the fan-out operator
triggers on structural consumer fan-out rather than on the composite score, as intended --- but the
per-edit acceptance test that should reject edits failing to clear $\Delta I >
\kappa\,\sigma_{\text{seed}}$ (\S\ref{sec:acceptance}) is not implemented, and applying the full
compiled policy unconditionally produces a mixed result (\S\ref{sec:mixed-result}): a positive mean
across scenarios that is not a reliable win, since three of seven scenarios come out worse. Closing
the loop from ``which component, and why'' to ``what to change, confirmed before deployment''
therefore remains partially open, and we report it as such rather than as an accomplished result.

Fourth, automated quality gating operationalises these checks continuously (RQ4). By implementing
in-memory evaluations using the \texttt{MemoryRepository} to bypass Neo4j database dependencies, the
framework executes anti-pattern scans and counterfactual failure simulations in seconds (\textasciitilde5~s
for medium, \textasciitilde40~s for xlarge). This execution speed allows the graph-based analyzer to
run as a blocking CI/CD build check, and the delta-aware gate semantics (\S\ref{sec:gate}) make the
check sustainable in practice: it blocks newly introduced architectural regressions at commit time
--- bridging the ``Architecture-Code Gap'' --- without repeatedly flagging known, risk-accepted
structure, analogous to the ``Clean as You Code'' discipline of static code analysis gates.

\subsection{Threats to Validity}

\textbf{Construct validity.} Our ground-truth impact is produced by a discrete-event simulator
rather than observed in deployed systems. The strongest claims we can make are therefore
comparative: our results speak to which modeling choices perform better under identical conditions,
not to absolute predictive accuracy in operation. High concordance with $I(v)$ indicates that a
predictor has captured the simulator's notion of cascade semantics; whether that notion matches a
specific production system is a separate, empirical question. We mitigate this by evaluating every
predictor against the same simulator-derived targets, so comparisons among the interpretable,
learned, homogeneous, and structural predictors remain internally consistent.

\textbf{Internal validity.} The chief internal risk is circular validation --- a predictor scoring
well because its inputs leaked from its labels. The framework's independence guarantee addresses
this directly: predictors operate on $G_{\text{analysis}}$ while ground truth is generated by
simulating $G_{\text{structural}}$, no simulation output is fed back as a predictor feature, and the
remediation stage generates candidates without reading $I(v)$. A measured correlation therefore
reflects predictive content rather than leakage. We further report bootstrap confidence intervals and
paired significance tests so that comparative claims are not artifacts of seed variance.

\textbf{External validity.} The synthetic suite, while spanning seven deployment domains and a range
of scales and structural regimes, is generated rather than harvested from production systems, and
this paper reports no validation against a real-world deployment or human expert judgment. We reduce
the in-distribution/out-of-distribution concern with Leave-One-Scenario-Out evaluation, which tests
transfer to held-out architectures and closes the transductive-leakage gap, but this evaluates
transfer across \emph{synthetic} scenarios only. Generalization to production systems, to other
middleware families, and to systems with richer adversarial structure remains future work
(\S\ref{sec:futurework}).

\subsection{Limitations and Future Work}
\label{sec:futurework}

Several limitations point to concrete next steps. \S\ref{sec:results}.1's in-distribution comparison
reports RMAV/$Q$ and HGL figures measured under different conventions (validated-dataset best case
versus an all-scenario mean); assembling a single per-scenario head-to-head table for RMAV/$Q$, HGL,
and $HGL\text{-}QoS$ on identical splits would pin down the exact margin rather than only its
direction. The Vulnerability dimension is the lightest of the four, resting on reachability-style
proxies; a richer adversarial model (trust boundaries, privilege escalation paths) would strengthen
the V attribution and broaden the framework's security relevance. This paper validates entirely on
synthetic scenarios; the most important next step is external validation on one or more real-world
pub-sub deployments, ideally including a blind expert-ranking study comparing predicted criticality
to domain-expert judgment, to test whether the ordering the framework recovers on synthetic
topologies agrees with practitioner assessment of an operational system. The remediation operator
set is small and deliberately conservative; expanding it, and deriving the acceptance multiplier
$\kappa$ from broader multi-seed variance data, would let the prescriptive stage address more failure
modes. Most concretely, \S\ref{sec:acceptance}'s per-edit acceptance test is not yet implemented as
an actual filter: \texttt{PrescribeService} applies the full compiled policy unconditionally rather
than simulating each candidate edit individually and discarding those that fail $\Delta I >
\kappa\,\sigma_{\text{seed}}$, and \S\ref{sec:mixed-result} shows this produces a mixed result
(positive mean, but negative in 3 of 7 scenarios). Implementing that filter and re-measuring is a
precondition for the remediation stage's practical-value claim to be more than aspirational, and is
the single most concrete next step this paper identifies. Finally, the entire framework is validated
against simulation; the natural endpoint is calibration against, or replacement of, the simulated
ground truth with observed failure data from instrumented deployments, which would convert the
comparative claims of this paper into absolute ones.

\subsection{Conclusion}

We presented Software-as-a-Graph, a pre-deployment Static System Analysis (SSA) framework that
models distributed pub-sub middleware as a typed, weighted, directed multigraph and analyzes it
along two coupled axes: multi-dimensional quality attribution, which decomposes each component's
criticality into orthogonal, interpretable RMAV dimensions (integrating local code quality metrics),
and failure-impact analysis, which predicts cascade impact with both the interpretable composite and
a learned heterogeneous graph transformer, validated against discrete-event simulation under a
strict input--label independence guarantee. A prescriptive remediation stage generates
topology-level hardening edits from structure alone and measures their effect end to end against the
canonical simulator; we report that measurement honestly as a mixed result, since the per-edit
acceptance test that would filter out counterproductive edits is not yet implemented
(\S\ref{sec:acceptance}, \S\ref{sec:mixed-result}).

Integrated directly into pipelines as a delta-aware, blocking CI/CD Quality Gate, the framework
verifies architectural changes and blocks regression in seconds, bridging the ``Architecture-Code
Gap'' at commit time. Across a synthetic scenario suite, the framework shows when interpretable
attribution suffices and when learning is required, exposes failure modes --- a shared-library blast
radius and a node-type stratification effect --- that single-score centrality cannot, and
demonstrates that attribution computed entirely before deployment can be made both legible and
actionable. By taking the \emph{type} of every component and dependency seriously, the framework
recovers structure that untyped, single-dimensional methods discard, and does so at the point in the
lifecycle where it is most valuable: before the system runs.

% ===========================================================================
\section*{References}
% ===========================================================================

\begin{thebibliography}{99}

\bibitem{Eugster2003} P. T. Eugster, P. A. Felber, R. Guerraoui, A.-M. Kermarrec, ``The many faces
of publish/subscribe,'' \emph{ACM Computing Surveys}, vol. 35, no. 2, pp. 114--131, 2003.

\bibitem{OMG2015DDS} Object Management Group, ``Data Distribution Service (DDS),'' OMG Document
formal/2015-04-10, version 1.4, 2015.

\bibitem{OASIS2019MQTT} OASIS, ``MQTT Version 5.0,'' OASIS Standard, 2019.

\bibitem{Freeman1977} L. C. Freeman, ``A set of measures of centrality based on betweenness,''
\emph{Sociometry}, vol. 40, no. 1, pp. 35--41, 1977.

\bibitem{BrinPage1998} S. Brin, L. Page, ``The anatomy of a large-scale hypertextual web search
engine,'' \emph{Computer Networks and ISDN Systems}, vol. 30, no. 1--7, pp. 107--117, 1998.

\bibitem{Buldyrev2010} S. V. Buldyrev, R. Parshani, G. Paul, H. E. Stanley, S. Havlin,
``Catastrophic cascade of failures in interdependent networks,'' \emph{Nature}, vol. 464, pp.
1025--1028, 2010.

\bibitem{Fan2020Nature} C. Fan, L. Zeng, Y. Sun, Y.-Y. Liu, ``Finding key players in complex
networks through deep reinforcement learning,'' \emph{Nature Machine Intelligence}, vol. 2, pp.
317--324, 2020.

\bibitem{Fan2019CIKM} C. Fan, L. Zeng, Y. Ding, M. Chen, Y. Sun, Z. Liu, ``Learning to identify high
betweenness centrality nodes from scratch: A novel graph neural network approach,'' in \emph{Proc.
28th ACM Int. Conf. on Information and Knowledge Management (CIKM)}, 2019, pp. 559--568.

\bibitem{Varbella2024PowerGraph} A. Varbella, K. Amara, M. El-Assady, B. Gjorgiev, G. Sansavini,
``PowerGraph: A power grid benchmark dataset for graph neural networks,'' in \emph{Advances in
Neural Information Processing Systems 37 (NeurIPS 2024), Datasets and Benchmarks Track}, 2024.
arXiv:2402.02827.

\bibitem{Schlichtkrull2018RGCN} M. Schlichtkrull, T. N. Kipf, P. Bloem, R. van den Berg, I. Titov, M.
Welling, ``Modeling relational data with graph convolutional networks,'' in \emph{Proc. European
Semantic Web Conference (ESWC)}, 2018, pp. 593--607.

\bibitem{Wang2019HAN} X. Wang, H. Ji, C. Shi, B. Wang, Y. Ye, P. Cui, P. S. Yu, ``Heterogeneous
graph attention network,'' in \emph{Proc. The Web Conference (WWW)}, 2019, pp. 2022--2032.

\bibitem{Hu2020HGT} Z. Hu, Y. Dong, K. Wang, Y. Sun, ``Heterogeneous graph transformer,'' in
\emph{Proc. The Web Conference (WWW)}, 2020, pp. 2704--2710.

\bibitem{Fu2020MAGNN} X. Fu, J. Zhang, Z. Meng, I. King, ``MAGNN: Metapath aggregated graph neural
network for heterogeneous graph embedding,'' in \emph{Proc. The Web Conference (WWW)}, 2020, pp.
2331--2341.

\bibitem{Li2018DeeperInsights} Q. Li, Z. Han, X.-M. Wu, ``Deeper insights into graph convolutional
networks for semi-supervised learning,'' in \emph{Proc. AAAI Conference on Artificial Intelligence},
2018, pp. 3538--3545.

\bibitem{Saaty1980} T. L. Saaty, \emph{The Analytic Hierarchy Process: Planning, Priority Setting,
Resource Allocation}, McGraw-Hill, 1980.

\bibitem{AnonA} Anonymized for review --- authors' prior work on multi-layer graph dependency
analysis for pub-sub architectures.

\end{thebibliography}

\end{document}
